% !TEX TS-program = pdflatex
% -----------------------------------------------------------------------------
% File: v17.tex
% Compile:
%   pdflatex v17.tex
%   pdflatex v17.tex
% -----------------------------------------------------------------------------
\documentclass[11pt]{article}

% --- Machine-readable PDF / font hygiene (pdfLaTeX) ---
\usepackage{cmap}
\usepackage[T1]{fontenc}
\usepackage[utf8]{inputenc}
\usepackage{lmodern}
\IfFileExists{glyphtounicode.tex}{\input{glyphtounicode.tex}}{}
\pdfgentounicode=1
\usepackage{microtype}
\usepackage{setspace}
\setstretch{1.2}

\usepackage[margin=0.7in]{geometry}
\setlength{\emergencystretch}{3em}
\setlength{\parskip}{0.35em}
\setlength{\parindent}{0em}
\sloppy

\usepackage{amsmath,amssymb,amsthm}
\usepackage{booktabs}
\usepackage{array}
\usepackage{enumitem}
\usepackage{xcolor}
\usepackage{hyperref}
\hypersetup{
  pdftitle={Observable-Only Proof-Carrying Autonomy (OOPCA): Audit Compression and Hybrid Proof/Replay Gating for No-Meta Agents},
  pdfauthor={K. Takahashi},
  pdfsubject={No-meta, observable-only autonomy, proof-carrying audit, hybrid gating, verification-limited scaling},
  pdfkeywords={no-meta, observable-only, proof-carrying data, SNARK, IVC, deterministic replay, fail-closed, audit compression, verification-limited scaling, accountability},
  colorlinks=true, linkcolor=black, citecolor=black, urlcolor=black
}

\usepackage{listings}
\lstset{
  basicstyle=\ttfamily\footnotesize,
  columns=fullflexible,
  breaklines=true,
  breakatwhitespace=false,
  breakautoindent=false,
  breakindent=0pt,
  postbreak=\mbox{\textcolor{gray}{$\hookrightarrow$}\space},
  frame=single,
  xleftmargin=1em,
  xrightmargin=1em
}

% ----------------------------
% Normative language (RFC 2119 / RFC 8174 style, used informally)
% ----------------------------
\newcommand{\MUST}{\textbf{MUST}}
\newcommand{\MUSTNOT}{\textbf{MUST NOT}}
\newcommand{\SHOULD}{\textbf{SHOULD}}
\newcommand{\SHOULDNOT}{\textbf{SHOULD NOT}}
\newcommand{\MAY}{\textbf{MAY}}

% ----------------------------
% Notation helpers
% ----------------------------
\newcommand{\Hash}{\mathsf{H}}
\newcommand{\Commit}{\mathsf{Commit}}
\newcommand{\Verify}{\mathsf{Verify}}
\newcommand{\Prove}{\mathsf{Prove}}
\newcommand{\Enc}{\mathsf{Enc}}
\newcommand{\Hex}{\mathsf{Hex}}
\newcommand{\Rel}{\mathcal{R}}
\newcommand{\Policy}{\mathcal{P}}
\newcommand{\Trace}{\mathcal{T}}
\newcommand{\Stmt}{\mathsf{S}}
\newcommand{\Proof}{\pi}

% ----------------------------
% Theorem environments
% ----------------------------
\newtheorem{definition}{Definition}
\newtheorem{lemma}{Lemma}
\newtheorem{theorem}{Theorem}
\newtheorem{claim}{Claim}
\newtheorem{corollary}{Corollary}

\title{Observable-Only Proof-Carrying Autonomy (OOPCA)\\[0.5em]
\large Audit Compression and Hybrid Proof/Replay Gating for No-Meta Agents}
\author{K. Takahashi\\\small\href{https://orcid.org/0009-0004-4273-3365}{ORCID 0009-0004-4273-3365}}
\date{February 2, 2026}

\begin{document}
\maketitle

\begin{abstract}
No-meta, observable-only autonomous agents cannot rely on privileged external judges and must ground safety-critical decisions in publicly observable evidence. As such systems scale, shared verification becomes a dominant bottleneck: verifiers must ingest, replay, and audit growing logs under adversarial participation.
This paper introduces \emph{Observable-Only Proof-Carrying Autonomy (OOPCA)}, a protocol-level audit mode that compresses verifier-side auditing by replacing parts of deterministic replay with succinct proof verification while preserving fail-closed semantics and non-inflationary progress accounting.
OOPCA is \emph{hybrid by default}: proofs are used only for policy components whose semantics are pinned to a deterministic \emph{Quantized Deterministic Semantics (QDS) profile} and compiled into a pinned relation artifact; deterministic replay remains the canonical fallback for components that are not safely compiled (e.g., floating-point dynamics, large unstructured search, and debugging/diagnostics).
OOPCA further (i) pins cryptographic and encoding conventions (hash algorithm, domain separation, and canonical JSON) to prevent verifier divergence, (ii) binds each gate decision to an explicit \emph{action intent} and an \emph{evidence manifest} so proofs cannot be reused for different actions or hidden inputs, (iii) provides positive cutoff certificates based on deterministic work-counter state updates (avoiding succinct ``absence'' claims), (iv) enforces bounded liveness using positive fallback-unlock receipts and an observable fallback ledger, and (v) supports accountability via proof-carrying minimal diagnostics plus Merkle-based controlled trace disclosure.
Finally, we provide a self-contained, machine-consumable schema bundle aligned with the ONCAP vocabulary.
\end{abstract}

\section{Introduction}

No-meta autonomous agents operate without a privileged external evaluator. In observable-only settings, safety-critical decisions must be derived from evidence that is publicly inspectable and deterministically checkable. The core design motif is \emph{fail-closed gating}: an action is permitted only when required evidence is present and verifies under pinned rules; otherwise the gate remains closed by default \cite{Takahashi2026ONCAP}.

A practical scaling bottleneck emerges. Replay-based audit systems require verifiers to ingest and re-execute large logs. Even when replay is deterministic, the \emph{shared} verification workload accumulates as systems grow, and the limiting resource becomes verification capacity itself \cite{Takahashi2026VLIA}. This suggests an audit substrate shift: from ``replay the whole log'' to ``verify a short proof''.

However, naïvely replacing replay with succinct proofs creates new failure modes: verifier divergence due to encoding ambiguity, semantic drift between CPU execution and finite-field circuits, the hardness of proving search cutoffs or ``best effort,'' liveness collapse when proofs cannot be generated, loss of diagnostics and accountability, reuse of proofs for different actions, and governance paradoxes around circuit/key updates.

OOPCA addresses these hazards directly. OOPCA does not claim that all auditing can be replaced by proofs. Instead, it defines a \emph{hybrid} audit mode that:
\begin{itemize}[leftmargin=1.2em]
  \item uses proofs only where constraints are explicitly quantized and circuit-friendly;
  \item uses replay receipts where compilation would introduce unacceptable drift or trusted-base expansion;
  \item preserves fail-closed integrity and non-inflationary progress definitions in both modes;
  \item binds gate decisions to explicit action intent and evidence manifests, preventing proof reuse and hidden inputs;
  \item provides verifiable, minimal diagnostics and controlled disclosure mechanisms for accountability;
  \item avoids trusted-clock assumptions by deriving windows from observable action turns.
\end{itemize}

\subsection{Claims and Non-Claims}

\textbf{Claims.} OOPCA provides:
\begin{itemize}[leftmargin=1.2em]
  \item verifier-side audit compression when proof mode is applicable;
  \item gate decisions based only on observable evidence (proof receipts and/or replay receipts) and explicitly declared evidence manifests;
  \item explicit conditions under which proof mode refines replay mode under QDS-bound semantics;
  \item positive cutoff certificates (work exhaustion) instead of succinct non-existence claims;
  \item bounded-liveness fallbacks with deterministic exposure bounds tracked by observable fallback ledgers;
  \item an explicit packet acceptance order that mitigates verification-capacity DoS.
\end{itemize}

\textbf{Non-claims.} OOPCA \emph{does not} claim:
\begin{itemize}[leftmargin=1.2em]
  \item that total computation decreases (cost often shifts from verifier to prover);
  \item that all policies can be compiled to efficient circuits;
  \item that software bugs or bad specifications vanish;
  \item that upgrades can be solved without trade-offs in no-meta settings.
\end{itemize}

\subsection{Threat Model}

We assume adversaries may:
\begin{itemize}[leftmargin=1.2em]
  \item deviate from intended execution while attempting to pass gates (including adaptive behavior conditioned on what verifiers accept);
  \item forge, malform, or replay packets; attempt cross-action reuse of previously valid receipts; or exploit ambiguity to induce verifier divergence;
  \item exploit weaknesses in proof systems (invalid proofs that nonetheless verify, malformed verifying keys, or soundness breaks);
  \item mount denial-of-service attacks by spamming verifiers with oversized, deeply nested, or expensive-to-parse objects, or by withholding required artifacts to stall progress;
  \item attempt governance attacks (key substitution, parameter poisoning, malicious upgrades, witness quorum capture, rollback/rewind attempts).
\end{itemize}

\paragraph{Observable-only constraint.}
Verifiers do not access private traces unless they are explicitly disclosed as additional observable evidence under a pinned rule (Section~\ref{sec:diag}).
All consensus-critical checks are functions of observable bytes only.

\paragraph{Cryptographic assumptions (explicit).}
OOPCA's safety statements rely on standard assumptions:
(i) collision resistance and preimage resistance of SHA-256 for all identifier hashes,
(ii) unforgeability of the signature scheme used for signer and witness receipts (any secure scheme may be used, but the scheme identifier \MUST\ be pinned by policy),
and (iii) soundness of the chosen proof system for the declared relation when in proof mode.
If (iii) is suspected to fail, OOPCA requires a bug-witness response that forces fail-closed closure for affected artifacts (Section~\ref{sec:setup}).

\paragraph{What is \emph{not} assumed.}
OOPCA does \emph{not} assume honest majority, a trusted clock, a trusted execution environment, or a privileged meta-evaluator.
Witness quorums are treated as \emph{observable cosigners}: they can improve liveness and accountability, but no safety claim depends on their honesty; compromise triggers fail-closed containment rather than silent acceptance.

\subsection{Roadmap}

Section~\ref{sec:related} positions OOPCA relative to proof-carrying approaches. Section~\ref{sec:prelim} sets notation and pins cryptographic/encoding conventions, QDS, and trace commitment profiles. Section~\ref{sec:core} defines statements, manifests, receipts, and hybrid gating. Section~\ref{sec:refinement} formalizes replay-to-proof refinement with an explicit specification-to-relation gap. Section~\ref{sec:cutoff} introduces positive cutoff certificates. Section~\ref{sec:liveness} provides bounded-liveness fallbacks. Section~\ref{sec:diag} covers diagnostics and controlled disclosure. Section~\ref{sec:setup} addresses witness quorums, key binding, and deterministic bug closure. Appendix~\ref{app:schemas} provides a self-contained JSON Schema bundle.

\section{Related Work and Positioning}\label{sec:related}

OOPCA inherits the observable-only, no-meta stance and fail-closed gating discipline from ONCAP \cite{Takahashi2026ONCAP}, and targets the scaling bottleneck characterized in verification-limited acceleration \cite{Takahashi2026VLIA}. It connects to floor and contract mechanisms for adversarial participation and intersubjective constraint enforcement \cite{Takahashi2025FSNA,Takahashi2025FMIC} and to causal-loop integrity constraints for long-horizon, self-referential operation \cite{Takahashi2025CLI}.

Outside this program, OOPCA is adjacent to:
\begin{itemize}[leftmargin=1.2em]
  \item proof-carrying approaches where code or data carry proofs of predicate satisfaction \cite{Necula1997PCC};
  \item succinct arguments (SNARK/STARK families) and incrementally verifiable computation (IVC) \cite{BenSasson2014SNARK,Groth2016,BenSasson2018STARK,Setty2021Nova};
  \item accumulation schemes and proof-carrying data (PCD) that support recursive composition and succinct verification across data structures \cite{Bunz2020PCD};
  \item zkVM and proof-at-the-boundary approaches.
\end{itemize}

\paragraph{What OOPCA adds.}
OOPCA emphasizes the following constraints simultaneously:
\begin{itemize}[leftmargin=1.2em]
  \item \textbf{Observable-only decision}: gate open/close is determined from a pinned, checkable packet set; no privileged judge.
  \item \textbf{Hybrid-by-default}: proof mode is primary where safe, but deterministic replay is a first-class fallback.
  \item \textbf{Pinned conventions}: hash algorithm, domain separation, and canonical encoding are fixed to prevent verifier divergence.
  \item \textbf{Action and evidence binding}: each receipt binds to an action intent and to an explicit evidence manifest (no hidden inputs, no proof reuse).
  \item \textbf{QDS as pinned artifact}: arithmetic and encoding rules are hash-fixed and checkable; an ASCII label alone is insufficient.
  \item \textbf{Merklized trace commitments}: controlled disclosure is possible without full trace disclosure, under a fully specified Merkle profile.
  \item \textbf{Cutoff as positive exhaustion}: proof of work-budget exhaustion under deterministic update rules, not proof of non-existence.
  \item \textbf{Positive-evidence liveness with bounded exposure}: fallback requires an explicit unlock receipt, and usage is tracked by verifiable receipts.
\end{itemize}

\section{Preliminaries}\label{sec:prelim}

\subsection{Turns, Epochs, and Windows}

OOPCA does not assume a trusted wall clock. Any ``time-like'' index used for budgeting or liveness is derived from observable evidence.
In this paper, a \emph{turn} $t\in\mathbb{N}$ is the action-sequence index \texttt{action\_seq\_u64} carried by the corresponding \texttt{ActionIntent} packet.
A \emph{window} is the closed interval $[t_0,t_1]$; by default OOPCA uses single-turn windows with $t_0=t_1=t$ for each action.

Verifiers \MUST\ reject any manifest or receipt whose \texttt{window\_start\_turn\_u64}/\texttt{window\_end\_turn\_u64} are inconsistent with the bound action intent's \texttt{action\_seq\_u64} under the policy's declared mapping.
Policies \MAY\ define coarser windows only by specifying an explicit, fully observable mapping rule; absent such a rule, the single-turn default applies.

\subsection{Pinned Cryptographic and Encoding Conventions}\label{sec:conventions}

\paragraph{Normative language.}
The key words \MUST, \MUSTNOT, \SHOULD, \SHOULDNOT, and \MAY\ are to be interpreted as described in RFC~2119 and RFC~8174 \cite{RFC2119,RFC8174}.

Observable-only systems fail in subtle ways if verifiers interpret the same bytes differently. OOPCA therefore pins minimal conventions that all verifiers must share.

\begin{definition}[Hash function selection and identifiers]\label{def:hash}
OOPCA uses a single cryptographic hash function $\Hash_{\mathrm{alg}}$ with 32-byte output, selected by the policy via \texttt{hash\_alg\_ascii\_id}.
A conformant verifier \MUST\ support \texttt{hash:sha256} (SHA-256 over raw bytes) \cite{FIPS1804}.
A policy MAY select other 32-byte-output algorithms only from an explicitly enumerated allowlist; such a selection is consensus-critical and therefore \emph{pinned} by \texttt{policy\_id\_hex} and \texttt{policy\_epoch\_u64}.
\end{definition}


\begin{definition}[Domain separation encoding]\label{def:ds}
For any domain string \texttt{ds} (ASCII lowercase, pattern \texttt{\string^[a-z0-9][a-z0-9:\_\-\.]{0,63}\string$}), define the domain-separated hash
\[
\Hash(\texttt{ds}, m) := \Hash_{\mathrm{alg}}\!\Big(\texttt{"OOPCA"} \Vert 0x00 \Vert \texttt{ASCII(hash\_alg\_ascii\_id)} \Vert 0x00 \Vert \Enc_{\text{U8}}(|\texttt{ds}|) \Vert \texttt{ASCII(ds)} \Vert 0x00 \Vert m\Big),
\]
where \texttt{hash\_alg\_ascii\_id} is selected by the policy (Definition~\ref{def:hash}) and therefore pinned by \texttt{policy\_id\_hex}.
Here $|\texttt{ds}|$ is the byte-length of \texttt{ASCII(ds)} (so $|\texttt{ds}|\le 64$), and $\Enc_{\text{U8}}(\cdot)$ is a single-byte length encoding.
\end{definition}

\paragraph{Domain allowlist.}
The policy field \texttt{allowed\_ds\_ascii\_ids} is a consensus-critical allowlist for domain strings.
A conformant verifier \MUST\ reject any policy that omits any domain string required to verify that policy's receipts and commitments.



\begin{definition}[Canonical JSON encoding]\label{def:jcs}
OOPCA uses a single, pinned canonical JSON encoding to prevent verifier divergence.
All consensus-critical objects are encoded using RFC~8785 JSON Canonicalization Scheme (JCS) over UTF-8 JSON bytes \cite{RFC8785}.
In addition to RFC~8785, the following consensus rules apply:
\begin{itemize}[leftmargin=1.2em]
\item \textbf{Duplicate keys are forbidden.} Parsers \MUSTNOT accept objects with duplicate member names at any nesting level.
\item \textbf{No NaN/Infinity.} JSON numbers \MUSTNOT represent NaN or infinities. Implementations \MUST\ reject any non-finite number encountered.
\item \textbf{Integer-only critical fields.} All consensus-critical counters and turns \MUST\ be encoded as decimal strings (schemas use \texttt{dec\_u64}, etc.).
\item \textbf{Set-like arrays are sorted and unique.} Any array declared as a set \MUST\ be strictly increasing under the comparator specified for that array and contain no duplicates. For arrays of strings, the comparator is lexicographic order on UTF-8 bytes. For \texttt{payload\_refs}, sort by the tuple \texttt{(payload\_id\_ascii\_id, payload\_hash\_hex)}. For \texttt{required\_families} and \texttt{family\_receipt\_bindings}, sort by \texttt{family\_ascii\_id}.
\item \textbf{Parse-cost caps are enforced before hashing.} Verifiers \MUST\ enforce policy caps on maximum nesting depth, maximum keys per object, maximum array length, maximum string bytes, and maximum packet bytes, \emph{before} computing any identifier hash or verifying any signature/proof.
\end{itemize}
\end{definition}

\begin{definition}[Base64 encoding (consensus-critical)]\label{def:b64}
All \texttt{*\_b64} fields are RFC~4648 base64 using the standard alphabet \texttt{A--Z a--z 0--9 + /} with padding \texttt{=} \cite{RFC4648}.
Whitespace and line breaks are forbidden.
Verifiers \MUST\ reject non-base64 strings and \MUST\ enforce a single canonical encoding (in particular: correct padding and no alternative alphabets).
\end{definition}
\begin{definition}[Packet envelope and packet identifier hash]\label{def:packetid}
OOPCA treats every observable artifact as a \emph{packet} with an explicit kind tag and a canonical JSON body.
Let a packet be an ordered pair \texttt{(packet\_kind\_ascii\_id, packet\_body)} where \texttt{packet\_kind\_ascii\_id} is an ASCII identifier and \texttt{packet\_body} is a JSON object encoded by Definition~\ref{def:jcs}.
Define the packet-ID preimage object \texttt{Env} as the JCS object with members \texttt{packet\_kind\_ascii\_id} and \texttt{packet\_body}.
The packet identifier hash is:
\[
\texttt{packet\_id\_hash\_hex} := \Hex(\Hash(\texttt{"oopca:packet-id"}, \Enc_{JCS}(\texttt{Env}))).
\]
Here $\Enc_{JCS}(\cdot)$ denotes RFC~8785 canonical JSON bytes, $\Hash$ is Definition~\ref{def:hash}, and $\Hex$ outputs lowercase hexadecimal.
Verifiers \MUST compute \texttt{packet\_id\_hash\_hex} only after enforcing policy parse-cost caps (Definition~\ref{def:jcs}).
\end{definition}
\subsection{Evidence Sets and Closure}\label{sec:evidence}
OOPCA verification is performed against an explicit \emph{evidence set} declared by an evidence-manifest packet and closed under pinned reference rules.

\begin{definition}[Evidence set]\label{def:evidence}
Fix an action intent packet and an evidence-manifest packet that references it.
The evidence set for a gate decision consists of:
\begin{itemize}[leftmargin=1.2em]
\item the action-intent packet (the target of the gate),
\item the evidence-manifest packet itself,
\item the policy packet identified by \texttt{policy\_id\_hex} (i.e., the policy packet's \texttt{packet\_id\_hash\_hex} per Definition~\ref{def:packetid}), and any policy-extension packets explicitly referenced by policy,
\item every packet whose identifier hash appears in \texttt{packet\_id\_hashes\_hex} of the evidence-manifest, and
\item every payload blob whose (id,hash) appears in \texttt{payload\_refs} across the manifest and included packets.
\end{itemize}
All elements of the evidence set are observable bytes. No hidden inputs are permitted.
\end{definition}

\paragraph{Closure rule.} 
Let \texttt{Allow} be the set of packet identifier hashes listed in \texttt{packet\_id\_hashes\_hex} together with the action intent, manifest, and policy packets.
For every included packet body, the verifier \MUST\ extract every referenced \texttt{*\_packet\_id\_hash} field (including nested occurrences, and including \texttt{OptionalHashRef} values with \texttt{mode=some}) by walking the JSON value in JCS key order and array order.
Each referenced hash \MUST\ be an element of \texttt{Allow}; otherwise the evidence set is malformed and verification \MUST\ fail closed.
Similarly, every referenced \texttt{payload\_refs} entry \MUST\ be present in the observable payload pool, and its hash \MUST\ match.

\paragraph{Deterministic extraction surface.} 
To make reference extraction deterministic across implementations, OOPCA treats a JSON string member whose key matches the suffix \texttt{\_packet\_id\_hash} and whose value matches \texttt{hex32} as a packet-id reference.
This surface is purely syntactic and is compatible with the schema bundle in Appendix~\ref{app:schemas}.

\begin{definition}[Evidence-manifest completeness]\label{def:manifestrule}
The evidence manifest is the observable ``closure certificate'' for an evidence set.
It \MUST\ include (as explicit fields) the identifiers of the policy packet and action-intent packet, and it \MUST\ enumerate (in \texttt{packet\_id\_hashes\_hex}) every additional packet required to evaluate the decision under the policy.
The following rules are consensus-critical:
\begin{itemize}[leftmargin=1.2em]
  \item \textbf{No self-inclusion}: \texttt{packet\_id\_hashes\_hex} \MUST\ NOT contain the identifier of the evidence-manifest packet itself.
  \item \textbf{Core packets are not listed}: \texttt{packet\_id\_hashes\_hex} \MUST\ NOT contain the identifiers of the policy packet or the action-intent packet (they are carried separately to avoid cycles and redundant bindings).
  \item \textbf{Sorted set}: \texttt{packet\_id\_hashes\_hex} is a set-like array and \MUST\ be strictly lexicographically sorted with no duplicates.
  \item \textbf{Closure}: every packet referenced by any included packet body (via a \texttt{\_packet\_id\_hash} field) \MUST\ either be one of the core packets or appear in \texttt{packet\_id\_hashes\_hex}. Likewise, every payload referenced by any included packet \MUST\ be present in the evidence set.
  \item \textbf{Receipt bindings are explicit:} the manifest \MUST\ include an explicit receipt-binding list (\texttt{family\_receipt\_bindings}) that identifies, per required family, the exact proof/replay and fallback receipts that are admissible for this decision. Receipts not bound by the manifest are non-admissible and cannot affect the gate result.
\end{itemize}
\end{definition}

This discipline makes ``hidden inputs'' mechanically hard: if a gate decision relies on some additional object, that object must appear as an observable packet or payload that is closed over by the manifest.
The cost is bounded by policy-defined budgets (e.g., maximum packet count and payload bytes), so DoS pressure is quantifiable and enforceable.


\subsection{Quantized Deterministic Semantics (QDS)}\label{sec:qds}

A critical obstacle to replay--proof alignment is arithmetic mismatch: real-world agent logic often uses floating point, while succinct proofs typically operate over finite fields. OOPCA resolves this by requiring that any behavior placed under proof mode be governed by an explicit \emph{Quantized Deterministic Semantics (QDS) profile}.

\begin{definition}[QDS profile artifact]
A QDS profile artifact is a hash-pinned object that specifies:
(i) numeric representation (integer-only or fixed-point with scale),
(ii) rounding/truncation conventions,
(iii) overflow behavior (wrap/saturate/error),
(iv) bit-width and signedness conventions for consensus-critical integers,
(v) canonical encoding for public inputs/outputs and trace commitments,
(vi) determinism constraints (no nondeterministic syscalls, randomized branches, or ambient timing),
(vii) admissible operations and forbidden sources of drift.
\end{definition}

\paragraph{Binding rule.}
A receipt \MUST\ bind to a QDS profile by including both:
\begin{itemize}[leftmargin=1.2em]
  \item a human-facing label \texttt{arith\_profile\_ascii\_id}, and
  \item a hash-pinned reference \texttt{arith\_profile\_packet\_id\_hash}.
\end{itemize}
Verifiers \MUST\ treat the QDS profile as the artifact referenced by \texttt{arith\_profile\_packet\_id\_hash}; labels are non-authoritative.

\paragraph{Consequence.}
Policy components not governed by a QDS profile artifact are \emph{ineligible} for proof mode and must be audited via deterministic replay receipts (or remain closed).

\subsection{Trace Commitment Profiles}\label{sec:tracecommit}

Controlled disclosure requires a commitment scheme that supports partial opening and does not leave tree-shape ambiguities to implementers.

\begin{definition}[Trace commitment profile]
A trace commitment profile is a pinned artifact specifying:
(i) canonical trace serialization (\texttt{trace\_serialization\_ascii\_id}),
(ii) leaf chunking (\texttt{chunk\_bytes\_u32}) and explicit length binding,
(iii) leaf hashing with domain separation,
(iv) Merkle tree construction, padding/odd-leaf rules, and root derivation \cite{Merkle1988},
(v) disclosure-path encoding (how a verifier reconstructs left/right ordering).
\end{definition}

\paragraph{Recommended construction (binary duplicate-last).}
Let the trace serialize to byte string $B$ with length $|B|$. Let $c=\texttt{chunk\_bytes\_u32}$ and $n=\lceil |B|/c \rceil$.
Define leaf $i$ bytes as $b_i=B[i\cdot c : \min((i+1)c,|B|)]$ with length $\ell_i$.
Define leaf hash:
\[
L_i := \Hash(\texttt{ds\_leaf}, \Enc_{\text{U32LE}}(i)\Vert \Enc_{\text{U32LE}}(\ell_i)\Vert b_i).
\]
Build the tree by pairing adjacent nodes and hashing internal nodes as:
\[
N := \Hash(\texttt{ds\_node}, L \Vert R),
\]
duplicating the last node when a level has odd count. The root is the final node.
Disclosure packets provide the sibling list \emph{and} direction bits so left/right ordering is unambiguous (Section~\ref{sec:disclosure}).
This construction is fully deterministic given $(B,c)$ and prevents ambiguity about last-chunk padding.

\subsection{Fail-Closed Gate Semantics}

\begin{definition}[Fail-closed gate]
A gate for action class $A$ is a function that outputs \texttt{OPEN} iff the required observable evidence verifies under pinned rules and the evidence manifest is complete. Otherwise it outputs \texttt{CLOSED}.
\end{definition}

Fail-closed prioritizes integrity over liveness; OOPCA therefore provides a bounded-liveness ladder (Section~\ref{sec:liveness}) rather than weakening this default.


\subsection{Verifier Determinism and Non-Divergence}\label{sec:determinism}

\begin{definition}[Conformant verifier]\label{def:conformant}
A verifier is \emph{conformant} if it:
(i) enforces all policy size caps before any expensive verification,
(ii) canonicalizes packet bodies using Definition~\ref{def:jcs},
(iii) recomputes every referenced packet identifier hash (all \texttt{*\_packet\_id\_hash} fields and any \texttt{OptionalHashRef.hash\_hex} with \texttt{mode=some}) using Definition~\ref{def:packetid},
(iv) evaluates gates using only the evidence set bound by the manifest (Section~\ref{sec:evidence}),
and (v) applies the hybrid gate logic in Appendix~\ref{app:pseudo} without reordering set-like arrays or relaxing rejection rules.
\end{definition}

\begin{theorem}[Gate determinism for a fixed observation frontier]\label{thm:determinism}
Fix a policy identifier \texttt{policy\_id\_hex}, epoch \texttt{policy\_epoch\_u64}, and a window $\langle \texttt{window\_start\_turn\_u64},\texttt{window\_end\_turn\_u64}\rangle$.
Assume the evidence manifest for that window \emph{explicitly binds} all consensus-critical anchors required for gating, including:
(i) the policy packet (hence \texttt{hash\_alg\_ascii\_id} and all pinned conventions), and
(ii) for any family that uses replay fallback in the manifest-bound receipt selection (Definition~\ref{def:manifest-bindings}), a fallback-ledger anchor under the policy's chosen ledger profile (\texttt{fallback\_ledger\_profile\_ascii\_id}).
In particular, under the GLOBAL ledger profile the manifest-bound evidence set \MUST\ include (and bind in \texttt{family\_receipt\_bindings}) a history-anchor checkpoint, a frontier checkpoint, and a consistency proof between them under a pinned \texttt{fallback\_ledger\_profile\_packet\_id\_hash} (Section~\ref{sec:fallback-ledger}).

Then any two conformant verifiers that observe the same manifest-bound evidence set output the same gate result (\texttt{OPEN} or \texttt{CLOSED}).
Different outputs are possible only when the inputs differ (e.g., different frontier checkpoints), in which case the ``freshness'' assumption is different and must be treated as an explicit operational parameter rather than hidden verifier-local state.
\end{theorem}

\begin{proof}
Under the assumptions, every acceptance predicate used by the verifier is a deterministic function of observable bytes:
the policy packet pins \texttt{hash\_alg\_ascii\_id} (Definition~\ref{def:hash}), domain separation (Definition~\ref{def:ds}), and canonical encoding (Definition~\ref{def:jcs});
the manifest fixes the admissible packet set and payload bytes; and the history-anchor packet makes any ledger-dependent checks a function of the evidence set rather than verifier-local memory.
Therefore, two conformant verifiers that observe the same evidence set compute the same digests and the same gate decision.
\end{proof}

\paragraph{Scope note (state divergence is not ``observable-only'').}
If a deployment chooses a \emph{verifier-local} fallback ledger without publishing a history anchor (or an equivalent globally checkable commitment), then different verifiers may legitimately diverge under network loss.
In that case, determinism holds only \emph{per verifier / per synchronized history}, and cross-verifier consensus is outside this paper's core protocol.


\section{OOPCA Core Model}\label{sec:core}

\subsection{Audit Modes}

OOPCA defines an audit mode for each gate family, declared by the policy field \texttt{mode}:
\begin{itemize}[leftmargin=1.2em]
  \item \textbf{\texttt{proof}}: the family is proof-mode. A proof receipt is required; a replay receipt is accepted only under an explicit fallback rule (positive unlock evidence plus bounded exposure accounting).
  \item \textbf{\texttt{replay}}: the family is replay-mode. A replay receipt is required.
\end{itemize}
The overall policy is \emph{hybrid} when it contains at least one proof-mode family and at least one replay-mode family; hybridness is not a third family mode.

\subsection{Action Intent Binding}

OOPCA binds each gate decision to an explicit action intent to prevent proof reuse.

\begin{definition}[Action intent]
An action intent is an observable packet describing the specific action request (class, parameters, and an action-scoped sequence) in canonical form. Gate evidence \MUST\ reference the action intent packet.
\end{definition}

\subsection{Statement Layer (Canonical Form)}\label{sec:statement}

A \emph{statement} is a canonical JSON object that fixes the public context of a gate decision for a single window. Importantly, the statement is a \emph{canonical object}, not a standalone packet kind; this avoids cyclic dependencies between receipts and manifests.

\begin{definition}[Statement canonical form]
Let \texttt{Stmt} be the canonical JSON object with the following fields:
\begin{itemize}[leftmargin=1.2em]
  \item \texttt{policy\_id\_hex}, \texttt{policy\_epoch\_u64},
  \item \texttt{window\_start\_turn\_u64}, \texttt{window\_end\_turn\_u64},
  \item \texttt{action\_class\_ascii\_id}, \texttt{action\_intent\_packet\_id\_hash},
  \item \texttt{arith\_profile\_ascii\_id}, \texttt{arith\_profile\_packet\_id\_hash},
  \item \texttt{trace\_commit\_root\_hash\_hex}, \texttt{trace\_commit\_profile\_packet\_id\_hash},
  \item \texttt{relation\_build\_receipt\_packet\_id\_hash},
  \item \texttt{diag\_pub}.
\end{itemize}
All values and nested objects \MUST\ match the schemas in Appendix~\ref{app:schemas}.
\end{definition}

\begin{definition}[Statement hash]
Let \texttt{statement\_hash\_hex} be:
\[
\texttt{statement\_hash\_hex} := \Hex(\Hash(\texttt{"oopca:statement-hash"}, \Enc_{\text{JCS}}(\texttt{Stmt}))).
\]
\end{definition}

\paragraph{Reconstruction rule.}
Verifiers \MUST\ reconstruct \texttt{Stmt} from receipt fields (and referenced packets) and recompute \texttt{statement\_hash\_hex}. Any mismatch \MUST\ be rejected.

\subsection{Evidence Manifest Binding}

The manifest binds an evidence set to a single window and policy epoch.

\begin{definition}[Manifest--statement binding]
The evidence manifest \MUST\ include \texttt{statement\_hash\_hex}. Every receipt in the same evidence set \MUST\ carry the same \texttt{statement\_hash\_hex}.
\end{definition}

This ``one hash'' binding avoids cycles: receipts do not reference the manifest identifier, and the manifest does not appear inside the statement.

\begin{definition}[Manifest receipt bindings]\label{def:manifest-bindings}
The evidence manifest \MUST\ include \path{family_receipt_bindings}, a set-like array with exactly one entry per policy-required family (\path{required_families[*].family_ascii_id}).
The policy's \path{required_families} list \MUST\ be sorted by \path{family_ascii_id} and contain no duplicate \path{family_ascii_id} values.
\path{family_receipt_bindings} \MUST\ be sorted by \path{family_ascii_id} and contain no duplicate \path{family_ascii_id} values.

Each entry carries \texttt{OptionalHashRef} pointers that select exactly one admissible receipt path for that family:
\begin{itemize}[leftmargin=1.2em]
  \item \textbf{Proof path (only if \texttt{mode=proof}):} \path{proof_receipt_packet_id_hash} is \texttt{mode=some}, \path{replay_receipt_packet_id_hash} is \texttt{mode=none}, and all fallback/ledger pointers are \texttt{mode=none}.
  \item \textbf{Replay path (only if \texttt{mode=replay}):} \path{replay_receipt_packet_id_hash} is \texttt{mode=some}, \path{proof_receipt_packet_id_hash} is \texttt{mode=none}, and all fallback/ledger pointers are \texttt{mode=none}.
  \item \textbf{Replay fallback (only if \texttt{mode=proof} and \texttt{fallback\_allowed\_bool=true}):} \path{proof_receipt_packet_id_hash} is \texttt{mode=none}; \path{replay_receipt_packet_id_hash} and \path{fallback_use_receipt_packet_id_hash} are \texttt{mode=some}; if \texttt{fallback\_unlock\_required\_bool=true} then \path{fallback_unlock_receipt_packet_id_hash} is also \texttt{mode=some}.
  The ledger-anchor pointers are mode-dependent: if \path{fallback_ledger_profile_ascii_id} is \path{global-log:v1}, the entry \MUST\ also bind \path{fallback_ledger_checkpoint_packet_id_hash}, \path{fallback_ledger_frontier_checkpoint_packet_id_hash}, and \path{fallback_ledger_consistency_proof_packet_id_hash} as \texttt{mode=some}; otherwise they \MUST\ be \texttt{mode=none}.
\end{itemize}
Any binding that sets both \texttt{proof\_receipt\_packet\_id\_hash} and \texttt{replay\_receipt\_packet\_id\_hash} to \texttt{mode=some} is rejected (determinism).
This makes receipt selection explicit and prevents adversarial ``receipt flooding'' from creating verifier divergence or liveness traps.
\end{definition}

\subsection{Proof Layer}

\begin{definition}[Proof receipt]
A proof receipt is an observable packet that includes:
(i) \texttt{statement\_hash\_hex} and the fields needed to reconstruct \texttt{Stmt},
(ii) a proof blob \texttt{proof\_b64} for the pinned relation, and
(iii) a cutoff-evidence reference \texttt{cutoff\_evidence\_packet\_id\_hash} (as \texttt{OptionalHashRef}; \texttt{mode=none} if unused by the family).
\end{definition}

\paragraph{Signature discipline.}
A proof receipt \MUST\ be authorized by (at least) one signer receipt targeting it (Section~\ref{sec:setup}). Proof receipts do not embed signer receipts to avoid cycles.

\subsection{Replay Layer}

\begin{definition}[Replay receipt]
A \emph{replay receipt} is an observable packet that certifies a gate evaluation via deterministic replay under a pinned replay profile.
It is the canonical fallback when proof mode is unavailable or unsafe (e.g., non-quantized dynamics, large unstructured search, or debugging builds).
A replay receipt \MUST\ include (at minimum):
\begin{itemize}[leftmargin=1.2em]
\item the window \texttt{(window\_start\_turn\_u64, window\_end\_turn\_u64)} and policy identifiers \texttt{(policy\_id\_hex, policy\_epoch\_u64)},
\item \texttt{action\_class\_ascii\_id} and \texttt{action\_intent\_packet\_id\_hash},
\item the pinned QDS profile identifiers \texttt{arith\_profile\_ascii\_id} and \texttt{arith\_profile\_packet\_id\_hash} (even if unused by a replay-mode family),
\item \texttt{trace\_commit\_root\_hash\_hex} and \texttt{trace\_commit\_profile\_packet\_id\_hash} binding the replay to an explicit trace-commit context,
\item \texttt{relation\_build\_receipt\_packet\_id\_hash} binding the evaluation to a specific compiled relation artifact when proofs are later used for the same statement,
\item \texttt{replay\_profile\_ascii\_id} and \texttt{replay\_profile\_packet\_id\_hash} identifying the deterministic replay profile,
\item \texttt{replay\_output\_digest\_hex}, a hash digest of the replay outputs relevant to the gate decision (including any bounded summaries required by policy), and
\item \texttt{diag\_pub}, a minimal public diagnostics object (Section~\ref{sec:diag}).
\end{itemize}
A verifier \MUST\ accept a replay receipt only if it (i) is included in the evidence set (Section~\ref{sec:evidence}), (ii) is bound to the same statement hash as any proof receipts for the same gate decision, and (iii) is consistent with the pinned replay profile and trace-commit rules.
\end{definition}

\paragraph{No ``hidden inputs'' claim.}
OOPCA does not attempt to prove the absence of private witness data inside a proof system. Instead, it requires that every dependency that can change a gate decision is either:
(i) included in \texttt{Stmt},
(ii) committed by \texttt{trace\_commit\_root\_hash\_hex}, or
(iii) an observable packet/payload enumerated by the manifest and checked under pinned conventions.

\subsection{Hybrid Gate Decision}

\begin{definition}[Hybrid gate open condition]
A gate is \texttt{OPEN} for a window iff all required families verify in their designated mode under the policy:
\begin{itemize}[leftmargin=1.2em]
  \item For a \textbf{proof-mode} family (\texttt{mode=proof}): a valid proof receipt is present, \emph{or} (only if \texttt{fallback\_allowed\_bool=true}) a valid replay receipt is present together with (if \texttt{fallback\_unlock\_required\_bool=true}) a valid fallback-unlock receipt and a valid fallback-use receipt (Section~\ref{sec:liveness}).
  \item For a \textbf{replay-mode} family (\texttt{mode=replay}): a valid replay receipt is present.
\end{itemize}
All receipts \MUST\ agree on \texttt{statement\_hash\_hex} and \texttt{trace\_commit\_root\_hash\_hex}.
\end{definition}

\subsection{Deterministic Acceptance Order}

To reduce DoS and eliminate non-determinism, verifiers \MUST\ apply the following order:
\begin{enumerate}[leftmargin=1.8em]
  \item Enforce policy caps (packet sizes, payload counts, and maximum receipt multiplicity).
  \item Parse the manifest; recompute its \texttt{packet\_id\_hash\_hex}; reject non-canonical bodies.
  \item Check manifest completeness (Section~\ref{sec:evidence}): all listed packets and all referenced payload bytes are present and hash-correct, and \texttt{family\_receipt\_bindings} refers only to included packets.
  \item Verify signer receipts for the action intent, all gate receipts, and all disclosure/fallback packets required by policy.
  \item Validate the relation build receipt under \texttt{build\_validation\_mode\_ascii\_id} (Section~\ref{sec:refinement}), and verify witness-quorum receipts only where the policy makes them consensus-critical (e.g., \texttt{trusted-witness:v1} build validation and bug witnesses).
  \item Verify replay receipts (deterministic replay), proof receipts (proof verification), and cutoff evidence as applicable.
  \item Apply the policy's family rules to decide \texttt{OPEN}/\texttt{CLOSED}.
\end{enumerate}

\paragraph{Mandatory cross-checks (\texttt{CrossChecksPass}).}
In addition to receipt-local verification, a conformant verifier \MUST\ enforce the following consensus-critical equalities and bindings across the manifest-bound evidence set; any mismatch \MUST\ fail closed:
\begin{itemize}[leftmargin=1.2em]
  \item \textbf{Policy binding:} all packets participating in the gate decision agree on \texttt{policy\_id\_hex} and \texttt{policy\_epoch\_u64}, and the policy packet pins \texttt{hash\_alg\_ascii\_id} and all convention identifiers used.
  \item \textbf{Window/action binding:} receipts bind to the manifest's \texttt{(window\_start\_turn\_u64, window\_end\_turn\_u64)} and \texttt{action\_intent\_packet\_id\_hash}, and these are consistent with the action intent's \texttt{action\_seq\_u64} under the policy's declared turn mapping.
  \item \textbf{Statement reconstruction:} \texttt{Stmt} reconstructed from the evidence set matches the receipt fields, and \texttt{statement\_hash\_hex} matches \texttt{Hex(Hash(..))} (Section~\ref{sec:statement}).
  \item \textbf{Trace context:} all selected receipts agree on \texttt{trace\_commit\_root\_hash\_hex} and \texttt{trace\_commit\_profile\_packet\_id\_hash}.
  \item \textbf{QDS binding:} all selected receipts agree on \texttt{arith\_profile\_ascii\_id} and \texttt{arith\_profile\_packet\_id\_hash}, and the referenced artifact is policy-admissible for any proof-eligible family.
  \item \textbf{Build/vk binding:} \texttt{relation\_build\_receipt\_packet\_id\_hash} referenced by receipts is present; its \texttt{vk\_hash\_hex} and \texttt{relation\_id\_hash\_hex} agree with receipt fields; and the build receipt validates under \texttt{build\_validation\_mode\_ascii\_id}.
  \item \textbf{Fallback ledger binding (when used):} if a proof-mode family opens via replay fallback, the verifier checks the manifest-bound fallback ledger anchor (LOCAL or GLOBAL) under the family's pinned \texttt{fallback\_ledger\_profile\_packet\_id\_hash} and enforces \texttt{fallback\_max\_windows\_u64} deterministically.
  \item \textbf{Reason digest binding:} \texttt{reason\_codes} (if present) canonicalize to a digest equal to \texttt{diag\_pub.reason\_digest\_hex} (Section~\ref{sec:diag}).
\end{itemize}
\section{Replay-to-Proof Refinement and the Specification-to-Relation Gap}\label{sec:refinement}

\subsection{Refinement profiles and the proof/replay acceptance gap}\label{sec:ref-profiles}

The proof relation $\Rel^{\text{proof}}$ and the replay relation $\Rel^{\text{replay}}$ are two executable semantics for the same policy-family logic.
From a security perspective, the safest target is \textbf{full equivalence} (accept/reject sets match).
From a systems perspective, full equivalence can be costly (e.g., arithmetic normalization, compilation, and circuit constraints).

OOPCA therefore makes the relationship \emph{explicit and policy-declared} via a \textbf{refinement profile} for each family (\texttt{refinement\_profile\_ascii\_id} in the family rule):

\begin{itemize}[leftmargin=1.2em]
  \item \textbf{EQUIVALENT}: $\Rel^{\text{proof}}(\Stmt,\Trace)=1 \iff \Rel^{\text{replay}}(\Stmt,\Trace)=1$ under pinned conventions (QDS/trace profiles). This profile is RECOMMENDED for safety- and liveness-critical families.
  \item \textbf{SOUND-SUBSET}: $\Rel^{\text{proof}}=1 \Rightarrow \Rel^{\text{replay}}=1$ only. This admits the possibility that proof is ``stricter'' than replay.
  \item \textbf{BOUNDED-GAP}: SOUND-SUBSET plus an explicit \texttt{gap\_budget\_u64} (policy-declared) that limits how many times a family may fall back to replay because proof did not materialize.
\end{itemize}

\paragraph{DoS and liveness degradation.}
If $\Rel^{\text{proof}}$ is a strict subset of $\Rel^{\text{replay}}$, an adversary can attempt to force repeated replay by choosing inputs that are valid under replay yet hard to prove.
OOPCA mitigates this by requiring that any such mismatch consumes scarce, auditable budget:
either the policy chooses EQUIVALENT, or it declares BOUNDED-GAP so that repeated ``proof-not-produced'' events are bounded and eventually fail-closed.
This aligns liveness with verification-limited scaling instead of allowing an attacker to push the system into perpetual replay.

\subsection{Pinned Relation Artifacts}

\begin{definition}[Relation build receipt]
A relation build receipt is an observable packet that binds:
\begin{itemize}[leftmargin=1.2em]
  \item \texttt{policy\_text\_hash\_hex} (policy source identifier),
  \item \texttt{arith\_profile\_packet\_id\_hash} (QDS artifact),
  \item \texttt{trace\_commit\_profile\_packet\_id\_hash},
  \item \texttt{relation\_id\_hash\_hex} (relation identifier),
  \item \texttt{vk\_hash\_hex} (verification key identifier),
  \item \texttt{toolchain\_hash\_hex} and reproducibility metadata,
  \item \texttt{build\_steps\_b64} (machine-checkable build script or digest thereof),
  \item an optional witness-set reference \texttt{witness\_set\_packet\_id\_hash} (as \texttt{OptionalHashRef}) when witnesses are used for coordination,
  \item and (optionally) a witness quorum receipt targeting the build receipt when the policy explicitly opts into trusted-witness build validation (Section~\ref{sec:setup}).
\end{itemize}
\end{definition}

\paragraph{Build-validation rule (no-meta by default).}
Verifiers \MUST\ treat a relation artifact as valid only if:
(i) hashes match pinned values and are consistent with the policy epoch's pinned identifiers,
(ii) the referenced QDS and trace-commit profile artifacts match the policy's declared proof-eligible components,
(iii) the pinned cryptographic and encoding conventions are satisfied (Section~\ref{sec:conventions}),
and
(iv) the build receipt is validated under the policy's declared \texttt{build\_validation\_mode\_ascii\_id}:
\begin{itemize}[leftmargin=1.2em]
  \item \textbf{\texttt{reproducible-build:v1}}: a verifier (or auditor) can reproduce the build from pinned inputs and confirm \texttt{relation\_id\_hash\_hex} and \texttt{vk\_hash\_hex}.
  \item \textbf{\texttt{proof-carrying-compilation:v1}}: the build includes a machine-checkable proof that the produced relation implements the pinned reference semantics under pinned conventions (e.g., translation validation or a verified compiler artifact) \cite{Necula2000TV,Leroy2009CompCert}.
  \item \textbf{\texttt{trusted-witness:v1}}: the build receipt is accepted if (and only if) a witness quorum receipt targeting it verifies under the policy-pinned witness set; this introduces a meta-level trust assumption and \MUST\ be stated in the deployment threat model.
\end{itemize}
Witness cosigns MAY be present in any mode as distribution/caching hints, but they are not sufficient for consensus-critical validity unless \texttt{trusted-witness:v1} is explicitly selected.

\subsection{Replay-to-Proof Refinement}

Let $\Rel^{\text{replay}}$ denote replay semantics under the same pinned QDS artifact for the proof-eligible kernel, and let $\Rel^{\text{proof}}$ denote the relation encoded by the pinned build receipt.

\begin{definition}[Refinement condition]
Proof mode refines replay mode for a component if, for all public inputs/outputs and traces under the pinned QDS artifact:
\[
\Rel^{\text{proof}}(\Stmt,\Trace)=1 \ \Rightarrow\ \Rel^{\text{replay}}(\Stmt,\Trace)=1.
\]
\end{definition}

This direction is safety-critical: if a proof is accepted, replay would also accept \emph{under the same QDS-bound semantics}. The reverse direction is optional and component-dependent.

\subsection{The Specification-to-Relation Gap}

Even with reproducible builds, equivalence between a high-level natural-language specification and a low-level relation is not automatic.

\begin{definition}[Spec-to-relation gap]
The spec-to-relation gap is the residual risk that the pinned relation artifact does not match the policy author's intent (bugs, miscompilation, ambiguous specs).
\end{definition}

OOPCA handles this gap by (i) narrowing the proof-eligible surface area to QDS-bound kernels, (ii) requiring explicit build receipts plus an explicit build-validation mode, and (iii) adopting a deterministic bug-closure regime (Section~\ref{sec:setup}) that prioritizes verifiable closure over unverifiable intervention.

\section{Cutoff Certificates and Non-Inflationary Progress Under Cutoffs}\label{sec:cutoff}

\subsection{Why ``absence proofs'' are the wrong target}

In many agent workloads, the hard part is not proving that a candidate satisfies constraints, but proving that \emph{no better} candidate exists or that ``best effort'' was expended. Such negative claims are generally not succinct without additional structure.

OOPCA therefore does \emph{not} attempt global non-existence proofs. Instead, it makes cutoffs \emph{positive and checkable}: the only claim proven is that a declared resource budget was consumed according to deterministic counter updates that are themselves part of the pinned relation (or replayed deterministically).

\paragraph{Limitation (Work-PoW is not Search-PoW).}
A cutoff certificate proves ``budget was spent under the pinned step rules''; it does \emph{not} prove that the budget was spent \emph{well} (e.g., that an optimizer explored a meaningful space, avoided redundant states, or approached an optimum).
Accordingly, OOPCA treats cutoff as a liveness/accounting artifact, not as a semantic guarantee of optimality or non-existence.
Policies that require stronger ``useful work'' properties must add additional structure (e.g., non-redundant exploration schedules, objective-improvement certificates, or challenge-seeded state transitions).

\subsection{Deterministic Work-Counter Model}\label{sec:countermodel}

Each proof-eligible search/optimization component declares a \texttt{counter\_model\_ascii\_id} and uses a deterministic state update that cannot be advanced ``for free'' without executing the pinned kernel.

A minimal counter model consists of:
\begin{itemize}[leftmargin=1.2em]
  \item a seed commitment \texttt{seed\_commit\_hash\_hex} bound to the statement hash,
  \item a step chain hash \texttt{step\_chain\_hash\_hex} updated per step using domain separation,
  \item a budget counter \texttt{budget\_used\_u64} incremented once per step.
\end{itemize}

\paragraph{Pinned update rule (example).}
Let $h_0 := \Hash(\texttt{"oopca:counter-step0"}, \texttt{seed\_commit} \Vert \texttt{statement\_hash})$. For step $i=0,1,\ldots$:
\[
h_{i+1} := \Hash(\texttt{"oopca:counter-step"}, h_i \Vert \Enc_{\text{U64LE}}(i) \Vert \Enc_{\text{JCS}}(\texttt{candidate}_i) \Vert \Enc_{\text{JCS}}(\texttt{eval}_i)).
\]
Here \texttt{candidate} is deterministically derived from the seed and index, and \texttt{eval} is the pinned kernel evaluation output under QDS. This prevents ``advancing the counter'' without producing QDS-consistent evaluation outputs.

\subsection{Cutoff Evidence}

\begin{definition}[Cutoff evidence]
Cutoff evidence for a window is an observable packet asserting \texttt{budget\_used\_u64 = budget\_cap\_u64} together with the deterministic counter model identifiers and terminal chain hash needed to check integrity. Cutoff evidence \MUST\ bind to the statement hash and trace-commit root.
\end{definition}

\paragraph{Key point.}
This is a \emph{positive} statement: the prover proves budget exhaustion under pinned update rules, not that no better solution exists.

\subsection{Non-Inflationary Progress Accounting}

Progress is defined strictly by \emph{verified evidence} and indexed by verified work, not by unverifiable claims of optimality.

\begin{definition}[Evidence-indexed progress credit]
For each window, define progress credit:
\[
\mathrm{Credit} := g(\texttt{policy\_id},\ \texttt{diag\_pub},\ \texttt{budget\_used})\cdot \mathbf{1}[\texttt{gate\_open}],
\]
where $g$ is a pinned deterministic function and \texttt{budget\_used} is verified (proof-verified for proof-eligible components; replay-verified otherwise).
\end{definition}

\begin{theorem}[No progress inflation under verified budgeting]\label{thm:npi}
Assume:
(i) the gate result for a window is computed by a conformant verifier from the evidence set bound by the manifest,
(ii) every term used in $\mathrm{Credit}$ is a deterministic function of that evidence set, and
(iii) any claim of work, optimality, or constraint satisfaction that is not supported by an accepted replay receipt, an accepted proof receipt, or an accepted cutoff evidence packet is ignored (fail-closed).
Then an adversary cannot increase $\mathrm{Credit}$ by introducing unverified statements, hidden evidence, or cross-action reuse of receipts.
\end{theorem}

\begin{proof}
Under (i), the verifier's admissible inputs are exactly the manifest-bound evidence set.
Under (ii), $\mathrm{Credit}$ is a deterministic function of those admissible inputs.
Under (iii), any additional unverified bytes are either rejected by syntax/canonicalization/manifest completeness checks or are simply non-admissible and therefore cannot change the accepted evidence set.
Cross-action reuse fails because receipts are bound to an action intent hash and statement hash (Section~\ref{sec:statement}).
Therefore the adversary cannot change the accepted evidence in a way that increases $\mathrm{Credit}$ without passing one of the verification predicates (replay, proof, or cutoff), which contradicts ``unverified'' by assumption.
\end{proof}

\section{Bounded Liveness Under Fail-Closed Gating}\label{sec:liveness}

Fail-closed gates can stall when a proof is not produced (e.g., infeasible compute, corrupted data, missing witnesses, or a miscompiled relation).
Under observable-only constraints it is invalid to open a gate because ``no proof appeared by a deadline'': absence is not observable evidence.
OOPCA therefore provides a \emph{positive-evidence} liveness ladder that preserves fail-closed safety while bounding how long an agent can remain blocked.

\subsection{Positive liveness ladder}

For a proof-mode family (\texttt{mode=proof}) with \texttt{fallback\_allowed\_bool=true}, a verifier MAY open the gate via exactly one of the following paths:

\begin{enumerate}
  \item \textbf{Proof path}: a valid \texttt{ProofReceipt} for the family and \texttt{statement\_hash\_hex}, authorized by the policy (signer receipts and, if required, witness quorum receipts).
  \item \textbf{Replay fallback path}: a valid \texttt{ReplayReceipt} \emph{plus} (if \texttt{fallback\_unlock\_required\_bool=true}) a valid \texttt{FallbackUnlockReceipt} authorizing fallback for the same family and \texttt{statement\_hash\_hex}, \emph{plus} a \texttt{FallbackUseReceipt} that advances the policy-defined fallback ledger without exceeding \texttt{fallback\_max\_windows\_u64}.
\end{enumerate}

No other path is permitted.
In particular, the replay fallback path never relies on absence claims; it requires a positive, auditable unlock artifact.

\subsection{Observable fallback ledger (ledger profiles and history anchors)}\label{sec:fallback-ledger}

Fallback is a liveness tool and therefore the most dangerous ``escape hatch'' in a fail-closed system.
OOPCA constrains fallback by requiring (i) \emph{positive authorization} (\texttt{FallbackUnlockReceipt}), and (ii) \emph{auditable bounded exposure} via a fallback ledger.

The review-critical point is that a \emph{verifier-local} ledger is not sufficient to guarantee \emph{cross-verifier} determinism under packet loss: different verifiers can miss different past fallback uses and therefore disagree about whether the exposure bound is exceeded.
To address this cleanly, OOPCA specifies \emph{ledger profiles}:

\begin{itemize}[leftmargin=1.2em]
  \item \textbf{LOCAL ledger profile (\texttt{fallback\_ledger\_profile\_ascii\_id}=\texttt{local:v1}) (minimal, not consensus-safe).}
  The verifier maintains ledger state derived from previously accepted evidence.
  This profile preserves bounded exposure \emph{for that verifier} but does \emph{not} guarantee cross-verifier agreement under network partitions.
  \item \textbf{GLOBAL ledger profile (\texttt{fallback\_ledger\_profile\_ascii\_id}=\texttt{global-log:v1}) (consensus-safe via published anchors).}
  Each fallback use is appended to a public, append-only transparency log (e.g., a Certificate Transparency-style log) or an equivalent globally checkable commitment stream \cite{RFC9162}.
  Verifiers require an inclusion proof (or checkpoint) so that exposure counts become a function of the evidence set, not local memory.
  The concrete checkpoint/proof verification rules are pinned by the family rule's \texttt{fallback\_ledger\_profile\_packet\_id\_hash}; labels are non-authoritative.
\end{itemize}

\paragraph{Ledger chain update (common core).}
Let $C_0$ be the fixed genesis chain hash defined by
\[
C_0 := \Hash(\texttt{"oopca:fallback-ledger:genesis"}, \Enc_{JCS}(\texttt{Env}_0)).
\]
where \texttt{Env}$_0$ is the JCS object with a single member \texttt{tag} whose value is the ASCII string \texttt{GENESIS}.
For each accepted fallback-use window, the observable fallback ledger advances deterministically as
\[
C_{k+1} := \Hash(\texttt{"oopca:fallback-ledger:update"}, \Enc_{JCS}(\texttt{Env}_{k+1})).
\]
Here \texttt{Env}$_{k+1}$ is the JCS object with members \texttt{prev\_chain\_hash\_hex} (equal to $\Hex(C_k)$), \texttt{statement\_hash\_hex}, and \texttt{family\_ascii\_id}.
A \texttt{FallbackUseReceipt} carries \texttt{fallback\_prev\_chain\_hash\_hex} (equal to $\Hex(C_k)$) and \texttt{fallback\_chain\_hash\_hex} (equal to $\Hex(C_{k+1})$); the verifier \MUST\ check that the update is consistent with the deterministic rule above.
The verifier \MUST\ also enforce \texttt{fallback\_count\_u64} $\le$ \texttt{fallback\_max\_windows\_u64} from the policy's family rule.

\paragraph{History anchors (required for cross-verifier determinism).}
In the GLOBAL profile, the evidence set \MUST\ include a \emph{history anchor} that commits to the current ledger head (and count) for the policy and family, together with a transparency-log inclusion proof or checkpoint chain.
In the LOCAL profile, the history anchor MAY be omitted; in that case, the exposure bound is an invariant only for a given verifier's accepted-history view.

\begin{definition}[Fallback-ledger history anchor packets]
In \texttt{global-log:v1}, the history anchor is represented by:
(i) a \texttt{FallbackLedgerCheckpoint} packet that commits to \texttt{(policy\_id\_hex, policy\_epoch\_u64, family\_ascii\_id, fallback\_count\_u64, fallback\_chain\_hash\_hex)} under a pinned \texttt{fallback\_ledger\_profile\_packet\_id\_hash}; and
(ii) a \texttt{FallbackLedgerConsistencyProof} packet that proves the claimed checkpoint is consistent with a selected \emph{frontier} checkpoint under the same pinned ledger profile.
These packet identifiers are bound explicitly in the manifest's \path{family_receipt_bindings} via \path{fallback_ledger_checkpoint_packet_id_hash}, \path{fallback_ledger_frontier_checkpoint_packet_id_hash}, and \path{fallback_ledger_consistency_proof_packet_id_hash}.
\end{definition}

\paragraph{Freshness, stale reads, and forks (no trusted clock).}
A transparency-log inclusion proof establishes that an anchor \emph{was} in a log, but in a no-meta setting it does not, by itself, establish that the anchor is \emph{the latest} head.
Without a trusted clock or a privileged ``latest-head oracle,'' a verifier cannot deduce freshness from a static evidence set: any past-valid checkpoint remains locally verifiable forever.
This is not a bug of OOPCA; it is the standard \emph{offline freshness impossibility} for global single-use constraints.

OOPCA therefore makes the required \emph{observation frontier} explicit.
In the GLOBAL ledger profile, the manifest-bound evidence set \MUST\ contain (i) a history-anchor checkpoint (e.g., a log head commitment) and (ii) a consistency proof linking that checkpoint to a \emph{frontier checkpoint} selected deterministically from the evidence set (Section~\ref{sec:fallback-ledger}).
Operationally, a conformant verifier acts as a \emph{monitor}: it maintains (or obtains via non-privileged gossip) a latest-seen frontier checkpoint for the policy and refuses to accept a checkpoint that cannot be proven consistent with that frontier.
If a verifier runs with a stale frontier, the global exposure bound holds only \emph{relative to that frontier}; this limitation is unavoidable without introducing a meta oracle.

\paragraph{Deterministic fork handling.}
If multiple frontier checkpoints are present in the evidence set, a conformant verifier \MUST\ select a unique frontier deterministically by the tuple
\[
(\texttt{checkpoint\_seq\_u64},\ \texttt{checkpoint\_hash\_hex})
\]
with lexicographic order on the hash as a tie-breaker, and then require a consistency proof from that frontier to the claimed checkpoint.
If no such proof is present, the verifier \MUST\ output \texttt{CLOSED}.

This yields the intended invariant: the replay fallback can restore progress, but only a bounded number of times, and only with positive authorization; and cross-verifier determinism is obtained exactly when the ledger head is made an observable, manifest-bound input.

\section{Diagnostics and Accountability Without Meta}\label{sec:diag}

Proof mode risks black-boxing: without replay logs, third parties may be unable to understand failures or debug bugs. OOPCA introduces \emph{proof-carrying diagnostics} and \emph{controlled disclosure} to restore accountability without requiring privileged judges.

\subsection{Two-Layer Diagnostics}

\textbf{Layer 1: minimal, verifiable diagnostics.}
A policy defines a minimal diagnostic object \texttt{diag\_pub} containing only circuit-friendly fields whose correctness is proven (or replay-verified). Examples:
\begin{itemize}[leftmargin=1.2em]
  \item \texttt{budget\_used\_u64}, \texttt{step\_count\_u64}, \texttt{step\_chain\_hash\_hex},
  \item \texttt{tri\_status} values (\texttt{sat}/\texttt{viol}/\texttt{uncert}),
  \item \texttt{closure\_class\_ascii\_id} (small classification),
  \item \texttt{reason\_digest\_hex} (hash binding for reason labels).
\end{itemize}

\textbf{Layer 2: optional reason labels.}
Packets include \texttt{reason\_codes} as an array of \texttt{ascii\_id} labels for interoperability and human inspection. These labels are non-authoritative unless bound by the digest below.

\subsection{Canonical Binding of Reason Labels}

To avoid expensive string handling inside circuits, OOPCA uses digest binding computed outside the circuit but checked against a proven (or replay-verified) digest.

\paragraph{Canonicalization rule.}
Define \texttt{reason\_codes\_canon} by:
\begin{itemize}[leftmargin=1.2em]
  \item ASCII lowercase normalization,
  \item rejection of non-ASCII or length $>64$,
  \item duplicate removal,
  \item lexicographic sort.
\end{itemize}
Then compute:
\[
\texttt{reason\_digest\_hex} := \Hex(\Hash(\texttt{ds\_reason}, \Enc_{\text{JCS}}(\texttt{reason\_codes\_canon}))).
\]
Verifiers \MUST\ compute this digest and compare to \texttt{diag\_pub.reason\_digest\_hex}. If the list is empty, the digest is the hash of the empty canonical list; therefore \texttt{reason\_digest\_hex} is always present.

\subsection{Controlled Disclosure: Merkle-Openable Trace Fragments}\label{sec:disclosure}

When stronger explanation is required, OOPCA allows optional disclosure of trace fragments as additional observable evidence:
\begin{itemize}[leftmargin=1.2em]
  \item a disclosure packet reveals bounded trace leaves plus Merkle openings to \texttt{trace\_commit\_root\_hash\_hex},
  \item the disclosure packet includes \texttt{path\_dirs\_b64} (direction bits) so left/right ordering is unambiguous,
  \item disclosure is policy-governed (bounded size, redaction rules, and rate limits).
\end{itemize}
Verifiers independently check the disclosed evidence against the pinned trace commitment profile.

\section{No-Meta Setup Hygiene, Witness Quorums, and Bug Response}\label{sec:setup}

OOPCA introduces relation build receipts and verification keys; without care, key updates can reintroduce meta-level authority.

\subsection{Witness Sets, quorum cosigning, and what they \emph{do not} guarantee}

Witness quorums are useful for coordination (publishing checkpoints, co-signing log roots, and accelerating distribution of known-good artifacts), but they are \emph{not} semantic oracles.
If a verifier accepts a relation build receipt \emph{solely} because a quorum cosigned it, then the system's safety claims implicitly depend on the honesty of that quorum---which is a meta-level trust assumption and must be stated as such.

To keep the ``no-meta'' constraint honest, OOPCA supports two deployment modes:

\begin{itemize}[leftmargin=1.2em]
  \item \textbf{Reproducible-build mode (no witness trust).}
  A build receipt \MUST\ include sufficient reproducibility metadata (toolchain hashes, source/input digests, deterministic build steps) such that any verifier (or any auditor) can independently reproduce the relation artifact and verify \texttt{relation\_id\_hash\_hex} and \texttt{vk\_hash\_hex}.
  In this mode, quorum cosigning is OPTIONAL and provides only availability / coordination value.
  \item \textbf{Proof-carrying compilation mode (minimal trusted base).}
  The build receipt includes a machine-checkable proof that the produced relation implements a stated reference semantics under pinned conventions (e.g., a proof about the compiler, a translation-validation proof, or a formally verified compiler artifact).
  In this mode, witnesses are again OPTIONAL; safety rests on the proof system and the pinned reference semantics.
\end{itemize}

\paragraph{Amortization and DoS hygiene for build validation.}
Reproducible builds are expensive, but they do not have to be on the per-window hot path.
OOPCA treats build validation as an \emph{epoch-bound} cost with explicit policy controls:
(i) epochs \MUST\ be long enough that recompilation is rare relative to per-window verification,
(ii) verifiers \MAY\ cache validated artifacts keyed by \texttt{relation\_artifact\_hash\_hex} (a pure function of inputs), and
(iii) policy \MUST\ bound the size and format of build inputs (toolchain metadata, source bundles, and flags) so that ``forced recompilation'' cannot be used as a denial-of-service lever.
When a verifier has not yet validated a new artifact, the safe default is fail-closed: \texttt{CLOSED} until validation completes.

\paragraph{Witness receipts as \emph{hints}, not authorities.}
A witness quorum can still be useful as a \emph{distribution and caching hint} (``many parties saw this artifact''), but under the no-meta threat model it is \emph{not} a substitute for either reproducible validation or a proof-carrying build.
This keeps the ``spec-to-relation'' link inside the observable evidence boundary rather than inside an unanalyzable trust assumption.


\paragraph{If a deployment chooses ``trusted witnesses''.}
A policy MAY still choose to treat the quorum as a trusted approval authority for pragmatic reasons.
In that case, the protocol is no longer strictly ``no-meta'' and the trust boundary must be explicit in the threat model.

\begin{definition}[Witness set declaration]
A witness set declaration is an observable packet that specifies:
(i) witness public keys (bytes) and their hashes,
(ii) a threshold,
(iii) the signature algorithm identifier,
(iv) deterministic rotation rules by epoch.

The arrays \texttt{witness\_pubkeys\_b64} and \texttt{witness\_pubkey\_hashes\_hex} \MUST\ have the same length and correspond index-wise.
\end{definition}

\begin{definition}[Witness quorum receipt]
A witness quorum receipt is an observable packet that binds:
(i) a target packet-ID hash (e.g., a build receipt packet),
(ii) the witness set packet-ID hash,
(iii) a threshold-valid signature bundle.
\end{definition}

\subsection{Pinned Key Binding}

Policies \MUST\ bind \texttt{policy\_id\_hex} and \texttt{policy\_epoch\_u64} to:
\begin{itemize}[leftmargin=1.2em]
  \item the policy source hash,
  \item the declared \texttt{build\_validation\_mode\_ascii\_id} for relation artifacts,
  \item (if any witness-quorum receipts are consensus-critical) the witness set packet-ID hash,
  \item the QDS profile packet-ID hash,
  \item the trace-commit profile packet-ID hash,
  \item the relation identifier hash,
  \item the verification key hash (and the verification key itself via payload reference).
\end{itemize}
Verifiers \MUST\ reject mismatches (fail-closed).

\subsection{Deterministic Rotation and Multi-Version Windows}

Policies \MAY\ include deterministic key-rotation schedules by epoch. To avoid abrupt liveness collapse, OOPCA allows a bounded transition window where both the old and new keys are accepted \emph{only if} each is bound to a pinned build receipt and each build receipt validates under \texttt{build\_validation\_mode\_ascii\_id}.
If \texttt{build\_validation\_mode\_ascii\_id} is \texttt{trusted-witness:v1}, each build receipt in the transition window \MUST\ also be cosigned by the active witness set under deterministic rotation rules.

\subsection{Deterministic Bug Closure}

If a relation bug is discovered, no-meta constraints prevent a privileged party from forcing an upgrade. OOPCA therefore adopts a deterministic bug-closure trigger.

\begin{definition}[Bug witness]
A bug witness is an observable packet, cosigned by the witness quorum (under the policy-pinned witness set), that identifies an affected build receipt (or relation/vk hash) and a closure scope (policy epoch range and/or window range). Upon observing a valid bug witness, verifiers deterministically close gates within scope.
\end{definition}

Correctness is preserved by closure; liveness recovery requires adopting a new epoch with new pinned artifacts under observable rules.

\section{Conclusion}

OOPCA defines an observable-only, no-meta audit mode that compresses verifier-side auditing by using succinct proofs where safe, while retaining deterministic replay as a first-class fallback. By pinning cryptographic and encoding conventions, binding each gate decision to explicit action intent and evidence manifests, pinning QDS and trace commitment profiles as artifacts (not labels), avoiding unrealistic full-equivalence claims via refinement, and replacing negative cutoff claims with positive deterministic work-exhaustion certificates, OOPCA preserves fail-closed integrity and non-inflationary progress accounting even under adversarial participation.

The protocol also restores accountability in proof mode via minimal verifiable diagnostics and optional Merkle-openable trace disclosures, and it maintains bounded liveness using positive fallback-unlock evidence with observable exposure accounting for replay fallback. Together these mechanisms shift shared verification effort away from replaying ever-growing logs and toward checking short, pinned receipts---without introducing a privileged meta-judge.

\appendix


\section{Self-Contained Schema Bundle (ONCAP-aligned)}\label{app:schemas}

This appendix provides a machine-consumable JSON Schema bundle that is self-contained (no external \texttt{\$ref}). It follows ONCAP naming discipline (e.g., \texttt{policy\_id\_hex}, \texttt{reason\_codes}, \texttt{OptionalHashRef}-style option wrappers) \cite{Takahashi2026ONCAP}.

\subsection{OOPCA Schema Bundle}


\paragraph{Validation rules beyond JSON Schema.}
The JSON Schema in this appendix constrains syntax, but several \emph{semantic} validation rules are consensus-critical and MUST be enforced by conformant verifiers:
(i) the packet-reference graph induced by \texttt{*\_packet\_id\_hash} fields (including nested occurrences and \texttt{OptionalHashRef} values with \texttt{mode=some}) MUST be acyclic (no cycles) with a fixed maximum reference depth;
(ii) packet references MUST be resolved under a fixed traversal order and total size cap;
(iii) any reference that would require unbounded recursion, repeated parsing, or repeated hashing is rejected (fail-closed);
(iv) \texttt{family\_receipt\_bindings} MUST satisfy Definition~\ref{def:manifest-bindings} (unique family coverage, mode-consistent receipt selection, and required ledger-anchor binding when fallback is used);
(v) all decimal-string integer fields \MUST\ be range-checked after parsing (reject overflow): \texttt{dec\_u16} in $[0,65535]$, \texttt{dec\_u32} in $[0,4294967295]$, \texttt{dec\_u64} in $[0,18446744073709551615]$, and \texttt{dec\_s32} in $[-2147483648,2147483647]$.
Payload blobs referenced by \texttt{payload\_refs} are treated as opaque bytes: verifiers do not recursively interpret payload bytes to discover further references.
These rules prevent parser- and resolver-level denial-of-service and make reference resolution deterministic.

\begin{lstlisting}
{
  "$schema": "https://json-schema.org/draft/2020-12/schema",
  "$id": "urn:oopca:schema:bundle:7",
  "title": "OOPCA Schema Bundle (v17)",
  "$defs": {
    "ascii_id": {
      "type": "string",
      "minLength": 1,
      "maxLength": 64,
      "pattern": "^[a-z0-9][a-z0-9:_\\-\\.]{0,63}$"
    },
    "hex32": {
      "type": "string",
      "pattern": "^[0-9a-f]{64}$"
    },
    "b64": {
      "type": "string",
      "minLength": 0,
      "maxLength": 1048576,
      "pattern": "^(?:[A-Za-z0-9+/]{4})*(?:[A-Za-z0-9+/]{2}==|[A-Za-z0-9+/]{3}=)?$"
    },
    "dec_u16": {
      "type": "string",
      "pattern": "^(0|[1-9][0-9]{0,4})$"
    },
    "dec_u32": {
      "type": "string",
      "pattern": "^(0|[1-9][0-9]{0,9})$"
    },
    "dec_u64": {
      "type": "string",
      "pattern": "^(0|[1-9][0-9]{0,19})$"
    },
    "dec_s32": {
      "type": "string",
      "pattern": "^(0|-?[1-9][0-9]{0,9})$"
    },
    "tri_status": {
      "type": "string",
      "enum": [
        "sat",
        "viol",
        "uncert"
      ]
    },
    "OptionalHashRef": {
      "type": "object",
      "required": [
        "mode"
      ],
      "properties": {
        "mode": {
          "type": "string",
          "enum": [
            "some",
            "none"
          ]
        },
        "hash_hex": {
          "$ref": "#/$defs/hex32"
        }
      },
      "additionalProperties": false,
      "allOf": [
        {
          "if": {
            "properties": {
              "mode": {
                "const": "some"
              }
            }
          },
          "then": {
            "required": [
              "hash_hex"
            ]
          }
        },
        {
          "if": {
            "properties": {
              "mode": {
                "const": "none"
              }
            }
          },
          "then": {
            "not": {
              "required": [
                "hash_hex"
              ]
            }
          }
        }
      ]
    },
    "PayloadRef": {
      "type": "object",
      "required": [
        "payload_id_ascii_id",
        "payload_hash_hex"
      ],
      "properties": {
        "payload_id_ascii_id": {
          "$ref": "#/$defs/ascii_id"
        },
        "payload_hash_hex": {
          "$ref": "#/$defs/hex32"
        }
      },
      "additionalProperties": false
    },
    "PayloadRefList": {
      "type": "array",
      "minItems": 0,
      "maxItems": 64,
      "items": {
        "$ref": "#/$defs/PayloadRef"
      },
      "uniqueItems": true
    },
    "OOPCAArithProfileBody": {
      "type": "object",
      "required": [
        "arith_profile_ascii_id",
        "hash_alg_ascii_id",
        "fixed_point_scale_s32",
        "word_bits_u16",
        "signed_bool",
        "round_mode_ascii_id",
        "overflow_mode_ascii_id",
        "encoding_ascii_id",
        "forbidden_sources",
        "payload_refs"
      ],
      "properties": {
        "arith_profile_ascii_id": {
          "$ref": "#/$defs/ascii_id"
        },
        "hash_alg_ascii_id": {
          "$ref": "#/$defs/ascii_id"
        },
        "fixed_point_scale_s32": {
          "$ref": "#/$defs/dec_s32"
        },
        "word_bits_u16": {
          "$ref": "#/$defs/dec_u16"
        },
        "signed_bool": {
          "type": "boolean"
        },
        "round_mode_ascii_id": {
          "$ref": "#/$defs/ascii_id"
        },
        "overflow_mode_ascii_id": {
          "$ref": "#/$defs/ascii_id"
        },
        "encoding_ascii_id": {
          "$ref": "#/$defs/ascii_id"
        },
        "forbidden_sources": {
          "type": "array",
          "minItems": 0,
          "maxItems": 64,
          "items": {
            "$ref": "#/$defs/ascii_id"
          },
          "uniqueItems": true
        },
        "payload_refs": {
          "type": "array",
          "minItems": 0,
          "maxItems": 64,
          "items": {
            "$ref": "#/$defs/PayloadRef"
          },
          "uniqueItems": true
        }
      },
      "additionalProperties": false
    },
    "OOPCATraceCommitProfileBody": {
      "type": "object",
      "required": [
        "trace_commit_profile_ascii_id",
        "hash_alg_ascii_id",
        "trace_serialization_ascii_id",
        "chunk_bytes_u32",
        "leaf_ds_ascii_id",
        "node_ds_ascii_id",
        "odd_leaf_rule_ascii_id",
        "path_encoding_ascii_id",
        "payload_refs"
      ],
      "properties": {
        "trace_commit_profile_ascii_id": {
          "$ref": "#/$defs/ascii_id"
        },
        "hash_alg_ascii_id": {
          "$ref": "#/$defs/ascii_id"
        },
        "trace_serialization_ascii_id": {
          "$ref": "#/$defs/ascii_id"
        },
        "chunk_bytes_u32": {
          "$ref": "#/$defs/dec_u32"
        },
        "leaf_ds_ascii_id": {
          "$ref": "#/$defs/ascii_id"
        },
        "node_ds_ascii_id": {
          "$ref": "#/$defs/ascii_id"
        },
        "odd_leaf_rule_ascii_id": {
          "$ref": "#/$defs/ascii_id"
        },
        "path_encoding_ascii_id": {
          "$ref": "#/$defs/ascii_id"
        },
        "payload_refs": {
          "type": "array",
          "minItems": 0,
          "maxItems": 64,
          "items": {
            "$ref": "#/$defs/PayloadRef"
          },
          "uniqueItems": true
        }
      },
      "additionalProperties": false
    },
    "OOPCAFallbackLedgerProfileBody": {
      "type": "object",
      "required": [
        "fallback_ledger_profile_ascii_id",
        "hash_alg_ascii_id",
        "payload_refs"
      ],
      "properties": {
        "fallback_ledger_profile_ascii_id": {
          "$ref": "#/$defs/ascii_id"
        },
        "hash_alg_ascii_id": {
          "$ref": "#/$defs/ascii_id"
        },
        "payload_refs": {
          "type": "array",
          "minItems": 0,
          "maxItems": 64,
          "items": {
            "$ref": "#/$defs/PayloadRef"
          },
          "uniqueItems": true
        }
      },
      "additionalProperties": false
    },
    "OOPCAFamilyReceiptBinding": {
      "type": "object",
      "required": [
        "family_ascii_id",
        "proof_receipt_packet_id_hash",
        "replay_receipt_packet_id_hash",
        "fallback_unlock_receipt_packet_id_hash",
        "fallback_use_receipt_packet_id_hash",
        "fallback_ledger_checkpoint_packet_id_hash",
        "fallback_ledger_frontier_checkpoint_packet_id_hash",
        "fallback_ledger_consistency_proof_packet_id_hash"
      ],
      "properties": {
        "family_ascii_id": {
          "$ref": "#/$defs/ascii_id"
        },
        "proof_receipt_packet_id_hash": {
          "$ref": "#/$defs/OptionalHashRef"
        },
        "replay_receipt_packet_id_hash": {
          "$ref": "#/$defs/OptionalHashRef"
        },
        "fallback_unlock_receipt_packet_id_hash": {
          "$ref": "#/$defs/OptionalHashRef"
        },
        "fallback_use_receipt_packet_id_hash": {
          "$ref": "#/$defs/OptionalHashRef"
        },
        "fallback_ledger_checkpoint_packet_id_hash": {
          "$ref": "#/$defs/OptionalHashRef"
        },
        "fallback_ledger_frontier_checkpoint_packet_id_hash": {
          "$ref": "#/$defs/OptionalHashRef"
        },
        "fallback_ledger_consistency_proof_packet_id_hash": {
          "$ref": "#/$defs/OptionalHashRef"
        }
      },
      "additionalProperties": false
    },
    "OOPCAEvidenceManifestBody": {
      "type": "object",
      "required": [
        "policy_id_hex",
        "policy_epoch_u64",
        "window_start_turn_u64",
        "window_end_turn_u64",
        "action_intent_packet_id_hash",
        "statement_hash_hex",
        "family_receipt_bindings",
        "packet_id_hashes_hex",
        "payload_refs"
      ],
      "properties": {
        "policy_id_hex": {
          "$ref": "#/$defs/hex32"
        },
        "policy_epoch_u64": {
          "$ref": "#/$defs/dec_u64"
        },
        "window_start_turn_u64": {
          "$ref": "#/$defs/dec_u64"
        },
        "window_end_turn_u64": {
          "$ref": "#/$defs/dec_u64"
        },
        "action_intent_packet_id_hash": {
          "$ref": "#/$defs/hex32"
        },
        "family_receipt_bindings": {
          "type": "array",
          "minItems": 1,
          "maxItems": 256,
          "items": {
            "$ref": "#/$defs/OOPCAFamilyReceiptBinding"
          },
          "uniqueItems": true
        },
        "packet_id_hashes_hex": {
          "type": "array",
          "minItems": 1,
          "maxItems": 4096,
          "items": {
            "$ref": "#/$defs/hex32"
          },
          "uniqueItems": true
        },
        "payload_refs": {
          "type": "array",
          "minItems": 0,
          "maxItems": 64,
          "items": {
            "$ref": "#/$defs/PayloadRef"
          },
          "uniqueItems": true
        },
        "statement_hash_hex": {
          "$ref": "#/$defs/hex32"
        }
      },
      "additionalProperties": false
    },
    "OOPCAActionIntentBody": {
      "type": "object",
      "required": [
        "policy_id_hex",
        "policy_epoch_u64",
        "action_class_ascii_id",
        "action_seq_u64",
        "action_params_b64",
        "payload_refs"
      ],
      "properties": {
        "policy_id_hex": {
          "$ref": "#/$defs/hex32"
        },
        "policy_epoch_u64": {
          "$ref": "#/$defs/dec_u64"
        },
        "action_class_ascii_id": {
          "$ref": "#/$defs/ascii_id"
        },
        "action_seq_u64": {
          "$ref": "#/$defs/dec_u64"
        },
        "action_params_b64": {
          "$ref": "#/$defs/b64"
        },
        "payload_refs": {
          "type": "array",
          "minItems": 0,
          "maxItems": 64,
          "items": {
            "$ref": "#/$defs/PayloadRef"
          },
          "uniqueItems": true
        }
      },
      "additionalProperties": false
    },
    "OOPCADiagnosticsPub": {
      "type": "object",
      "required": [
        "step_count_u64",
        "budget_used_u64",
        "step_chain_hash_hex",
        "tri_status",
        "closure_class_ascii_id",
        "reason_digest_hex"
      ],
      "properties": {
        "step_count_u64": {
          "$ref": "#/$defs/dec_u64"
        },
        "budget_used_u64": {
          "$ref": "#/$defs/dec_u64"
        },
        "step_chain_hash_hex": {
          "$ref": "#/$defs/hex32"
        },
        "tri_status": {
          "$ref": "#/$defs/tri_status"
        },
        "closure_class_ascii_id": {
          "$ref": "#/$defs/ascii_id"
        },
        "reason_digest_hex": {
          "$ref": "#/$defs/hex32"
        }
      },
      "additionalProperties": false
    },
    "OOPCAProofReceiptBody": {
      "type": "object",
      "required": [
        "action_class_ascii_id",
        "action_intent_packet_id_hash",
        "arith_profile_ascii_id",
        "arith_profile_packet_id_hash",
        "cutoff_evidence_packet_id_hash",
        "diag_pub",
        "payload_refs",
        "policy_epoch_u64",
        "policy_id_hex",
        "proof_b64",
        "reason_codes",
        "relation_build_receipt_packet_id_hash",
        "statement_hash_hex",
        "trace_commit_profile_packet_id_hash",
        "trace_commit_root_hash_hex",
        "vk_hash_hex",
        "window_end_turn_u64",
        "window_start_turn_u64"
      ],
      "properties": {
        "window_start_turn_u64": {
          "$ref": "#/$defs/dec_u64"
        },
        "window_end_turn_u64": {
          "$ref": "#/$defs/dec_u64"
        },
        "policy_id_hex": {
          "$ref": "#/$defs/hex32"
        },
        "policy_epoch_u64": {
          "$ref": "#/$defs/dec_u64"
        },
        "action_class_ascii_id": {
          "$ref": "#/$defs/ascii_id"
        },
        "action_intent_packet_id_hash": {
          "$ref": "#/$defs/hex32"
        },
        "arith_profile_ascii_id": {
          "$ref": "#/$defs/ascii_id"
        },
        "arith_profile_packet_id_hash": {
          "$ref": "#/$defs/hex32"
        },
        "relation_build_receipt_packet_id_hash": {
          "$ref": "#/$defs/hex32"
        },
        "statement_hash_hex": {
          "$ref": "#/$defs/hex32"
        },
        "trace_commit_root_hash_hex": {
          "$ref": "#/$defs/hex32"
        },
        "trace_commit_profile_packet_id_hash": {
          "$ref": "#/$defs/hex32"
        },
        "vk_hash_hex": {
          "$ref": "#/$defs/hex32"
        },
        "proof_b64": {
          "$ref": "#/$defs/b64"
        },
        "diag_pub": {
          "$ref": "#/$defs/OOPCADiagnosticsPub"
        },
        "reason_codes": {
          "type": "array",
          "minItems": 0,
          "maxItems": 64,
          "items": {
            "$ref": "#/$defs/ascii_id"
          },
          "uniqueItems": true
        },
        "cutoff_evidence_packet_id_hash": {
          "$ref": "#/$defs/OptionalHashRef"
        },
        "payload_refs": {
          "type": "array",
          "minItems": 0,
          "maxItems": 64,
          "items": {
            "$ref": "#/$defs/PayloadRef"
          },
          "uniqueItems": true
        }
      },
      "additionalProperties": false
    },
    "OOPCAReplayReceiptBody": {
      "type": "object",
      "required": [
        "action_class_ascii_id",
        "action_intent_packet_id_hash",
        "arith_profile_ascii_id",
        "arith_profile_packet_id_hash",
        "relation_build_receipt_packet_id_hash",
        "diag_pub",
        "payload_refs",
        "policy_epoch_u64",
        "policy_id_hex",
        "reason_codes",
        "replay_output_digest_hex",
        "replay_profile_ascii_id",
        "replay_profile_packet_id_hash",
        "statement_hash_hex",
        "trace_commit_profile_packet_id_hash",
        "trace_commit_root_hash_hex",
        "window_end_turn_u64",
        "window_start_turn_u64"
      ],
      "properties": {
        "window_start_turn_u64": {
          "$ref": "#/$defs/dec_u64"
        },
        "window_end_turn_u64": {
          "$ref": "#/$defs/dec_u64"
        },
        "policy_id_hex": {
          "$ref": "#/$defs/hex32"
        },
        "policy_epoch_u64": {
          "$ref": "#/$defs/dec_u64"
        },
        "action_class_ascii_id": {
          "$ref": "#/$defs/ascii_id"
        },
        "action_intent_packet_id_hash": {
          "$ref": "#/$defs/hex32"
        },
        "arith_profile_ascii_id": {
          "$ref": "#/$defs/ascii_id"
        },
        "arith_profile_packet_id_hash": {
          "$ref": "#/$defs/hex32"
        },
        "relation_build_receipt_packet_id_hash": {
          "$ref": "#/$defs/hex32"
        },
        "statement_hash_hex": {
          "$ref": "#/$defs/hex32"
        },
        "trace_commit_root_hash_hex": {
          "$ref": "#/$defs/hex32"
        },
        "trace_commit_profile_packet_id_hash": {
          "$ref": "#/$defs/hex32"
        },
        "replay_profile_ascii_id": {
          "$ref": "#/$defs/ascii_id"
        },
        "replay_profile_packet_id_hash": {
          "$ref": "#/$defs/hex32"
        },
        "replay_output_digest_hex": {
          "$ref": "#/$defs/hex32"
        },
        "diag_pub": {
          "$ref": "#/$defs/OOPCADiagnosticsPub"
        },
        "reason_codes": {
          "type": "array",
          "minItems": 0,
          "maxItems": 64,
          "items": {
            "$ref": "#/$defs/ascii_id"
          },
          "uniqueItems": true
        },
        "payload_refs": {
          "type": "array",
          "minItems": 0,
          "maxItems": 64,
          "items": {
            "$ref": "#/$defs/PayloadRef"
          },
          "uniqueItems": true
        }
      },
      "additionalProperties": false
    },
    "OOPCACutoffEvidenceBody": {
      "type": "object",
      "required": [
        "window_start_turn_u64",
        "window_end_turn_u64",
        "policy_id_hex",
        "policy_epoch_u64",
        "statement_hash_hex",
        "trace_commit_root_hash_hex",
        "budget_cap_u64",
        "budget_used_u64",
        "counter_model_ascii_id",
        "seed_commit_hash_hex",
        "step_chain_hash_hex",
        "payload_refs"
      ],
      "properties": {
        "window_start_turn_u64": {
          "$ref": "#/$defs/dec_u64"
        },
        "window_end_turn_u64": {
          "$ref": "#/$defs/dec_u64"
        },
        "policy_id_hex": {
          "$ref": "#/$defs/hex32"
        },
        "policy_epoch_u64": {
          "$ref": "#/$defs/dec_u64"
        },
        "statement_hash_hex": {
          "$ref": "#/$defs/hex32"
        },
        "trace_commit_root_hash_hex": {
          "$ref": "#/$defs/hex32"
        },
        "budget_cap_u64": {
          "$ref": "#/$defs/dec_u64"
        },
        "budget_used_u64": {
          "$ref": "#/$defs/dec_u64"
        },
        "counter_model_ascii_id": {
          "$ref": "#/$defs/ascii_id"
        },
        "seed_commit_hash_hex": {
          "$ref": "#/$defs/hex32"
        },
        "step_chain_hash_hex": {
          "$ref": "#/$defs/hex32"
        },
        "payload_refs": {
          "type": "array",
          "minItems": 0,
          "maxItems": 64,
          "items": {
            "$ref": "#/$defs/PayloadRef"
          },
          "uniqueItems": true
        }
      },
      "additionalProperties": false
    },
    "OOPCARelationBuildReceiptBody": {
      "type": "object",
      "required": [
        "arith_profile_packet_id_hash",
        "build_steps_b64",
        "compiler_ascii_id",
        "payload_refs",
        "policy_epoch_u64",
        "policy_id_hex",
        "policy_text_hash_hex",
        "relation_id_hash_hex",
        "toolchain_hash_hex",
        "trace_commit_profile_packet_id_hash",
        "vk_hash_hex",
        "witness_set_packet_id_hash"
      ],
      "properties": {
        "policy_id_hex": {
          "$ref": "#/$defs/hex32"
        },
        "policy_epoch_u64": {
          "$ref": "#/$defs/dec_u64"
        },
        "policy_text_hash_hex": {
          "$ref": "#/$defs/hex32"
        },
        "arith_profile_packet_id_hash": {
          "$ref": "#/$defs/hex32"
        },
        "trace_commit_profile_packet_id_hash": {
          "$ref": "#/$defs/hex32"
        },
        "compiler_ascii_id": {
          "$ref": "#/$defs/ascii_id"
        },
        "toolchain_hash_hex": {
          "$ref": "#/$defs/hex32"
        },
        "relation_id_hash_hex": {
          "$ref": "#/$defs/hex32"
        },
        "vk_hash_hex": {
          "$ref": "#/$defs/hex32"
        },
        "build_steps_b64": {
          "$ref": "#/$defs/b64"
        },
        "witness_set_packet_id_hash": {
          "$ref": "#/$defs/OptionalHashRef"
        },
        "payload_refs": {
          "type": "array",
          "minItems": 0,
          "maxItems": 64,
          "items": {
            "$ref": "#/$defs/PayloadRef"
          },
          "uniqueItems": true
        }
      },
      "additionalProperties": false
    },
    "OOPCAWitnessSetDeclareBody": {
      "type": "object",
      "required": [
        "policy_id_hex",
        "policy_epoch_u64",
        "sig_alg_ascii_id",
        "threshold_u16",
        "witness_pubkey_hashes_hex",
        "witness_pubkeys_b64",
        "rotation_rule_ascii_id",
        "payload_refs"
      ],
      "properties": {
        "policy_id_hex": {
          "$ref": "#/$defs/hex32"
        },
        "policy_epoch_u64": {
          "$ref": "#/$defs/dec_u64"
        },
        "sig_alg_ascii_id": {
          "$ref": "#/$defs/ascii_id"
        },
        "threshold_u16": {
          "$ref": "#/$defs/dec_u16"
        },
        "witness_pubkey_hashes_hex": {
          "type": "array",
          "minItems": 1,
          "maxItems": 128,
          "items": {
            "$ref": "#/$defs/hex32"
          }
        },
        "witness_pubkeys_b64": {
          "type": "array",
          "minItems": 1,
          "maxItems": 128,
          "items": {
            "$ref": "#/$defs/b64"
          }
        },

        "rotation_rule_ascii_id": {
          "$ref": "#/$defs/ascii_id"
        },
        "payload_refs": {
          "type": "array",
          "minItems": 0,
          "maxItems": 64,
          "items": {
            "$ref": "#/$defs/PayloadRef"
          },
          "uniqueItems": true
        }
      },
      "additionalProperties": false
    },
    "OOPCAWitnessQuorumReceiptBody": {
      "type": "object",
      "required": [
        "cosignatures",
        "payload_refs",
        "policy_epoch_u64",
        "policy_id_hex",
        "sig_alg_ascii_id",
        "target_packet_id_hash",
        "witness_set_packet_id_hash"
      ],
      "properties": {
        "policy_id_hex": {
          "$ref": "#/$defs/hex32"
        },
        "policy_epoch_u64": {
          "$ref": "#/$defs/dec_u64"
        },
        "witness_set_packet_id_hash": {
          "$ref": "#/$defs/hex32"
        },
        "target_packet_id_hash": {
          "$ref": "#/$defs/hex32"
        },
        "sig_alg_ascii_id": {
          "$ref": "#/$defs/ascii_id"
        },
        "payload_refs": {
          "type": "array",
          "minItems": 0,
          "maxItems": 64,
          "items": {
            "$ref": "#/$defs/PayloadRef"
          },
          "uniqueItems": true
        },
        "cosignatures": {
          "type": "array",
          "minItems": 1,
          "items": {
            "$ref": "#/$defs/OOPCACosignature"
          }
        }
      },
      "additionalProperties": false
    },
    "OOPCADisclosureBody": {
      "type": "object",
      "required": [
        "window_start_turn_u64",
        "window_end_turn_u64",
        "policy_id_hex",
        "policy_epoch_u64",
        "statement_hash_hex",
        "trace_commit_root_hash_hex",
        "trace_commit_profile_packet_id_hash",
        "leaf_index_u32",
        "leaf_len_u32",
        "leaf_bytes_b64",
        "path_dirs_b64",
        "merkle_siblings_hex",
        "payload_refs"
      ],
      "properties": {
        "window_start_turn_u64": {
          "$ref": "#/$defs/dec_u64"
        },
        "window_end_turn_u64": {
          "$ref": "#/$defs/dec_u64"
        },
        "policy_id_hex": {
          "$ref": "#/$defs/hex32"
        },
        "policy_epoch_u64": {
          "$ref": "#/$defs/dec_u64"
        },
        "statement_hash_hex": {
          "$ref": "#/$defs/hex32"
        },
        "trace_commit_root_hash_hex": {
          "$ref": "#/$defs/hex32"
        },
        "trace_commit_profile_packet_id_hash": {
          "$ref": "#/$defs/hex32"
        },
        "leaf_index_u32": {
          "$ref": "#/$defs/dec_u32"
        },
        "leaf_len_u32": {
          "$ref": "#/$defs/dec_u32"
        },
        "leaf_bytes_b64": {
          "$ref": "#/$defs/b64"
        },
        "path_dirs_b64": {
          "$ref": "#/$defs/b64"
        },
        "merkle_siblings_hex": {
          "type": "array",
          "minItems": 1,
          "maxItems": 128,
          "items": {
            "$ref": "#/$defs/hex32"
          }
        },
        "payload_refs": {
          "type": "array",
          "minItems": 0,
          "maxItems": 64,
          "items": {
            "$ref": "#/$defs/PayloadRef"
          },
          "uniqueItems": true
        }
      },
      "additionalProperties": false
    },
    "OOPCAFallbackUseReceiptBody": {
      "type": "object",
      "required": [
        "fallback_chain_hash_hex",
        "fallback_count_u64",
        "fallback_prev_chain_hash_hex",
        "family_ascii_id",
        "payload_refs",
        "policy_epoch_u64",
        "policy_id_hex",
        "statement_hash_hex",
        "window_end_turn_u64",
        "window_start_turn_u64"
      ],
      "properties": {
        "policy_id_hex": {
          "$ref": "#/$defs/hex32"
        },
        "policy_epoch_u64": {
          "$ref": "#/$defs/dec_u64"
        },
        "family_ascii_id": {
          "$ref": "#/$defs/ascii_id"
        },
        "window_start_turn_u64": {
          "$ref": "#/$defs/dec_u64"
        },
        "window_end_turn_u64": {
          "$ref": "#/$defs/dec_u64"
        },
        "statement_hash_hex": {
          "$ref": "#/$defs/hex32"
        },
        "fallback_prev_chain_hash_hex": {
          "$ref": "#/$defs/hex32"
        },
        "fallback_chain_hash_hex": {
          "$ref": "#/$defs/hex32"
        },
        "fallback_count_u64": {
          "$ref": "#/$defs/dec_u64"
        },
        "payload_refs": {
          "type": "array",
          "minItems": 0,
          "maxItems": 64,
          "items": {
            "$ref": "#/$defs/PayloadRef"
          },
          "uniqueItems": true
        }
      },
      "additionalProperties": false
    },
    "OOPCAFallbackLedgerCheckpointBody": {
      "type": "object",
      "required": [
        "checkpoint_hash_hex",
        "checkpoint_seq_u64",
        "fallback_chain_hash_hex",
        "fallback_count_u64",
        "family_ascii_id",
        "fallback_ledger_profile_ascii_id",
        "payload_refs",
        "policy_epoch_u64",
        "policy_id_hex"
      ],
      "properties": {
        "policy_id_hex": {
          "$ref": "#/$defs/hex32"
        },
        "policy_epoch_u64": {
          "$ref": "#/$defs/dec_u64"
        },
        "family_ascii_id": {
          "$ref": "#/$defs/ascii_id"
        },
        "fallback_ledger_profile_ascii_id": {
          "$ref": "#/$defs/ascii_id"
        },
        "checkpoint_seq_u64": {
          "$ref": "#/$defs/dec_u64"
        },
        "checkpoint_hash_hex": {
          "$ref": "#/$defs/hex32"
        },
        "fallback_count_u64": {
          "$ref": "#/$defs/dec_u64"
        },
        "fallback_chain_hash_hex": {
          "$ref": "#/$defs/hex32"
        },
        "payload_refs": {
          "type": "array",
          "minItems": 0,
          "maxItems": 64,
          "items": {
            "$ref": "#/$defs/PayloadRef"
          },
          "uniqueItems": true
        }
      },
      "additionalProperties": false
    },
    "OOPCAFallbackLedgerConsistencyProofBody": {
      "type": "object",
      "required": [
        "family_ascii_id",
        "fallback_ledger_profile_ascii_id",
        "from_checkpoint_hash_hex",
        "from_checkpoint_seq_u64",
        "payload_refs",
        "policy_epoch_u64",
        "policy_id_hex",
        "proof_b64",
        "to_checkpoint_hash_hex",
        "to_checkpoint_seq_u64"
      ],
      "properties": {
        "policy_id_hex": {
          "$ref": "#/$defs/hex32"
        },
        "policy_epoch_u64": {
          "$ref": "#/$defs/dec_u64"
        },
        "family_ascii_id": {
          "$ref": "#/$defs/ascii_id"
        },
        "fallback_ledger_profile_ascii_id": {
          "$ref": "#/$defs/ascii_id"
        },
        "from_checkpoint_seq_u64": {
          "$ref": "#/$defs/dec_u64"
        },
        "from_checkpoint_hash_hex": {
          "$ref": "#/$defs/hex32"
        },
        "to_checkpoint_seq_u64": {
          "$ref": "#/$defs/dec_u64"
        },
        "to_checkpoint_hash_hex": {
          "$ref": "#/$defs/hex32"
        },
        "proof_b64": {
          "$ref": "#/$defs/b64"
        },
        "payload_refs": {
          "type": "array",
          "minItems": 0,
          "maxItems": 64,
          "items": {
            "$ref": "#/$defs/PayloadRef"
          },
          "uniqueItems": true
        }
      },
      "additionalProperties": false
    },
    "OOPCABugWitnessBody": {
      "type": "object",
      "required": [
        "affected_build_receipt_packet_id_hash",
        "closure_scope_ascii_id",
        "payload_refs",
        "policy_epoch_u64",
        "policy_id_hex",
        "scope_window_end_turn_u64",
        "scope_window_start_turn_u64",
        "witness_set_packet_id_hash"
      ],
      "properties": {
        "policy_id_hex": {
          "$ref": "#/$defs/hex32"
        },
        "policy_epoch_u64": {
          "$ref": "#/$defs/dec_u64"
        },
        "affected_build_receipt_packet_id_hash": {
          "$ref": "#/$defs/hex32"
        },
        "closure_scope_ascii_id": {
          "$ref": "#/$defs/ascii_id"
        },
        "scope_window_start_turn_u64": {
          "$ref": "#/$defs/dec_u64"
        },
        "scope_window_end_turn_u64": {
          "$ref": "#/$defs/dec_u64"
        },
        "witness_set_packet_id_hash": {
          "$ref": "#/$defs/hex32"
        },
        "payload_refs": {
          "type": "array",
          "minItems": 0,
          "maxItems": 64,
          "items": {
            "$ref": "#/$defs/PayloadRef"
          },
          "uniqueItems": true
        }
      },
      "additionalProperties": false
    },
    "OOPCAPacketIdPreimage": {
      "type": "object",
      "required": [
        "packet_kind_ascii_id",
        "packet_body"
      ],
      "properties": {
        "packet_kind_ascii_id": {
          "$ref": "#/$defs/ascii_id"
        },
        "packet_body": {
          "type": "object"
        }
      },
      "additionalProperties": false
    },
    "OOPCAFamilyMode": {
      "type": "string",
      "enum": [
        "proof",
        "replay"
      ]
    },
    "OOPCABuildValidationMode": {
      "type": "string",
      "enum": [
        "reproducible-build:v1",
        "proof-carrying-compilation:v1",
        "trusted-witness:v1"
      ]
    },
    "OOPCAFamilyRule": {
      "type": "object",
      "required": [
        "fallback_allowed_bool",
        "fallback_ledger_profile_ascii_id",
        "fallback_ledger_profile_packet_id_hash",
        "fallback_max_windows_u64",
        "fallback_mode_ascii_id",
        "fallback_penalty_ascii_id",
        "fallback_unlock_required_bool",
        "family_ascii_id",
        "mode"
      ],
      "properties": {
        "family_ascii_id": {
          "$ref": "#/$defs/ascii_id"
        },
        "mode": {
          "$ref": "#/$defs/OOPCAFamilyMode"
        },
        "refinement_profile_ascii_id": {
          "$ref": "#/$defs/ascii_id"
        },
        "gap_budget_u64": {
          "$ref": "#/$defs/dec_u64"
        },
        "fallback_allowed_bool": {
          "type": "boolean"
        },
        "fallback_max_windows_u64": {
          "$ref": "#/$defs/dec_u64"
        },
        "fallback_unlock_required_bool": {
          "type": "boolean"
        },
        "fallback_mode_ascii_id": {
          "$ref": "#/$defs/ascii_id"
        },
        "fallback_penalty_ascii_id": {
          "$ref": "#/$defs/ascii_id"
        },
        "fallback_ledger_profile_ascii_id": {
          "$ref": "#/$defs/ascii_id"
        },
        "fallback_ledger_profile_packet_id_hash": {
          "$ref": "#/$defs/OptionalHashRef"
        }
      },
      "additionalProperties": false
    },
    "OOPCAPolicyBody": {
      "type": "object",
      "required": [
        "allowed_ds_ascii_ids",
        "authorized_signers",
        "build_validation_mode_ascii_id",
        "hash_alg_ascii_id",
        "max_evidence_packets_u32",
        "max_json_array_len_u32",
        "max_json_keys_u32",
        "max_json_nesting_depth_u16",
        "max_json_string_bytes_u32",
        "max_packet_bytes_u32",
        "max_payload_bytes_u32",
        "policy_epoch_u64",
        "required_families",
        "turn_mapping_ascii_id",
        "witness_set_packet_id_hash"
      ],
      "properties": {
        "policy_epoch_u64": {
          "$ref": "#/$defs/dec_u64"
        },
        "build_validation_mode_ascii_id": {
          "$ref": "#/$defs/OOPCABuildValidationMode"
        },
        "hash_alg_ascii_id": {
          "$ref": "#/$defs/ascii_id"
        },
        "max_packet_bytes_u32": {
          "$ref": "#/$defs/dec_u32"
        },
        "max_payload_bytes_u32": {
          "$ref": "#/$defs/dec_u32"
        },
        "max_evidence_packets_u32": {
          "$ref": "#/$defs/dec_u32"
        },
        "max_json_nesting_depth_u16": {
          "$ref": "#/$defs/dec_u16"
        },
        "max_json_keys_u32": {
          "$ref": "#/$defs/dec_u32"
        },
        "max_json_array_len_u32": {
          "$ref": "#/$defs/dec_u32"
        },
        "max_json_string_bytes_u32": {
          "$ref": "#/$defs/dec_u32"
        },

        "allowed_ds_ascii_ids": {
          "type": "array",
          "minItems": 1,
          "maxItems": 128,
          "items": {
            "$ref": "#/$defs/ascii_id"
          },
          "uniqueItems": true
        },
        "required_families": {
          "type": "array",
          "minItems": 1,
          "maxItems": 256,
          "items": {
            "$ref": "#/$defs/OOPCAFamilyRule"
          },
          "uniqueItems": true
        },
        "turn_mapping_ascii_id": {
          "enum": [
            "action_seq_u64"
          ]
        },
        "witness_set_packet_id_hash": {
          "$ref": "#/$defs/OptionalHashRef"
        },
        "authorized_signers": {
          "$ref": "#/$defs/OOPCASignerSetPolicy"
        }
      },
      "additionalProperties": false
    },
    "OOPCACosignature": {
      "type": "object",
      "additionalProperties": false,
      "required": [
        "witness_pubkey_hash_hex",
        "sig_alg_ascii_id",
        "sig_b64"
      ],
      "properties": {
        "witness_pubkey_hash_hex": {
          "$ref": "#/$defs/hex32"
        },
        "sig_alg_ascii_id": {
          "$ref": "#/$defs/ascii_id"
        },
        "sig_b64": {
          "$ref": "#/$defs/b64"
        }
      }
    },
    "OOPCASignerSetPolicy": {
      "type": "object",
      "additionalProperties": false,
      "required": [
        "sig_alg_ascii_id",
        "threshold_u16",
        "signer_pubkey_hashes_hex"
      ],
      "properties": {
        "sig_alg_ascii_id": {
          "$ref": "#/$defs/ascii_id"
        },
        "threshold_u16": {
          "$ref": "#/$defs/dec_u16"
        },
        "signer_pubkey_hashes_hex": {
          "type": "array",
          "minItems": 1,
          "items": {
            "$ref": "#/$defs/hex32"
          },
          "uniqueItems": true
        }
      }
    },
    "OOPCASignerReceiptBody": {
      "type": "object",
      "additionalProperties": false,
      "required": [
        "policy_id_hex",
        "policy_epoch_u64",
        "signer_pubkey_hash_hex",
        "sig_alg_ascii_id",
        "target_packet_id_hash",
        "signature_b64"
      ],
      "properties": {
        "policy_id_hex": {
          "$ref": "#/$defs/hex32"
        },
        "policy_epoch_u64": {
          "$ref": "#/$defs/dec_u64"
        },
        "signer_pubkey_hash_hex": {
          "$ref": "#/$defs/hex32"
        },
        "sig_alg_ascii_id": {
          "$ref": "#/$defs/ascii_id"
        },
        "target_packet_id_hash": {
          "$ref": "#/$defs/hex32"
        },
        "signature_b64": {
          "$ref": "#/$defs/b64"
        },
        "payload_refs": {
          "$ref": "#/$defs/PayloadRefList"
        }
      }
    },
    "OOPCAFallbackUnlockReceiptBody": {
      "type": "object",
      "additionalProperties": false,
      "required": [
        "policy_id_hex",
        "policy_epoch_u64",
        "window_start_turn_u64",
        "window_end_turn_u64",
        "action_intent_packet_id_hash",
        "statement_hash_hex",
        "family_ascii_id",
        "reason_code_ascii_id",
        "reason_digest_hex"
      ],
      "properties": {
        "policy_id_hex": {
          "$ref": "#/$defs/hex32"
        },
        "policy_epoch_u64": {
          "$ref": "#/$defs/dec_u64"
        },
        "window_start_turn_u64": {
          "$ref": "#/$defs/dec_u64"
        },
        "window_end_turn_u64": {
          "$ref": "#/$defs/dec_u64"
        },
        "action_intent_packet_id_hash": {
          "$ref": "#/$defs/hex32"
        },
        "arith_profile_ascii_id": {
          "$ref": "#/$defs/ascii_id"
        },
        "arith_profile_packet_id_hash": {
          "$ref": "#/$defs/hex32"
        },
        "relation_build_receipt_packet_id_hash": {
          "$ref": "#/$defs/hex32"
        },
        "statement_hash_hex": {
          "$ref": "#/$defs/hex32"
        },
        "family_ascii_id": {
          "$ref": "#/$defs/ascii_id"
        },
        "reason_code_ascii_id": {
          "$ref": "#/$defs/ascii_id"
        },
        "reason_digest_hex": {
          "$ref": "#/$defs/hex32"
        },
        "payload_refs": {
          "$ref": "#/$defs/PayloadRefList"
        }
      }
    }
  }
}
\end{lstlisting}

\subsection{Reason Code Enumeration (Recommended)}

Reason codes are \texttt{ascii\_id} labels. The following set is recommended for OOPCA mode:

\begin{center}
\begin{tabular}{@{}p{0.40\linewidth}p{0.52\linewidth}@{}}
\toprule
\textbf{reason\_code} & \textbf{meaning} \\ 
\midrule
\texttt{oopca:manifest\_missing} & evidence manifest required but absent \\ 
\texttt{oopca:manifest\_incomplete} & required packet not listed or missing \\ 
\texttt{oopca:action\_intent\_mismatch} & receipt does not bind to required action intent \\ 
\texttt{oopca:proof\_invalid} & proof receipt observed but verification failed \\ 
\texttt{oopca:proof\_missing} & required proof receipt absent and fallback path not satisfied \\ 
\texttt{oopca:fallback\_unlock\_missing} & fallback allowed but unlock receipt absent \\ 
\texttt{oopca:fallback\_exposure\_exceeded} & fallback ledger exceeds pinned exposure bounds \\ 
\texttt{oopca:replay\_fallback\_used} & gate opened via replay fallback under bounded rule \\ 
\texttt{oopca:qds\_profile\_mismatch} & QDS/arith profile artifact mismatch with policy requirement \\ 
\texttt{oopca:trace\_commit\_profile\_mismatch} & trace commitment profile mismatch \\ 
\texttt{oopca:build\_receipt\_unverified} & build receipt missing or quorum-cosign invalid \\ 
\texttt{oopca:cutoff\_exhausted} & cutoff evidence indicates verification budget exhausted \\ 
\texttt{oopca:bug\_observed} & bug witness evidence indicates invalidity \\ 
\texttt{oopca:diagnostics\_digest\_mismatch} & reason digest does not match canonical reason list \\ 
\texttt{oopca:disclosure\_path\_invalid} & Merkle opening/path encoding invalid \\ 
\texttt{oopca:disclosure\_submitted} & a trace-fragment disclosure was submitted \\ 
\bottomrule
\end{tabular}
\end{center}

\section{Pseudocode (Minimal)}\label{app:pseudo}

\subsection*{Verifier Gate Logic (Hybrid)}

\begin{lstlisting}
function VerifyGateHybrid(required_families, evidence_set):
  # Fail-closed default

  # Cheap prefilter (DoS discipline)
  if not SyntaxAndSizeCapsOK(evidence_set): return CLOSED
  if not CanonicalizationOK(evidence_set): return CLOSED

  # Manifest-first binding
  if not ManifestPresentAndValid(evidence_set): return CLOSED
  if not ManifestIsComplete(evidence_set): return CLOSED

  # Hash/action/trace bindings
  if not StatementHashOK(evidence_set): return CLOSED
  if not ActionIntentBindingOK(evidence_set): return CLOSED
  if not TraceCommitBindingOK(evidence_set): return CLOSED

  # Cosigns
  if not CosignaturesOK(evidence_set): return CLOSED

  for fam in required_families:
    if fam.mode == "proof":
      if not exists(evidence_set.proof_receipt[fam]):
        if fam.fallback_allowed_bool and exists(evidence_set.replay_receipt[fam]) and exists(evidence_set.fallback_use_receipt[fam]) and (not fam.fallback_unlock_required_bool or exists(evidence_set.fallback_unlock_receipt[fam])):
          if not VerifyReplayReceipt(evidence_set.replay_receipt[fam]): return CLOSED
          if fam.fallback_unlock_required_bool and not VerifyFallbackUnlockReceipt(evidence_set.fallback_unlock_receipt[fam]): return CLOSED
          if not VerifyFallbackUseReceipt(evidence_set.fallback_use_receipt[fam]): return CLOSED
          if not FallbackLedgerOK(fam, evidence_set): return CLOSED
          continue  # bounded fallback accepted
        return CLOSED
      if not VerifyProofReceipt(evidence_set.proof_receipt[fam]):
        return CLOSED
    else if fam.mode == "replay":
      if not exists(evidence_set.replay_receipt[fam]): return CLOSED
      if not VerifyReplayReceipt(evidence_set.replay_receipt[fam]): return CLOSED

  if not CrossChecksPass(evidence_set): return CLOSED
  return OPEN
\end{lstlisting}

\bibliographystyle{plain}
\begin{thebibliography}{30}

\bibitem{Takahashi2026ONCAP}
K.~Takahashi.
\newblock Observable-Only No-Meta Causal Autonomy Protocol (ONCAP).
\newblock Zenodo, 2026.
\newblock doi:10.5281/zenodo.18371930.

\bibitem{Takahashi2026VLIA}
K.~Takahashi.
\newblock Verification-Limited Intelligence Acceleration: Observable-Only Laws, Bounded Derivation, and Diagnostics under No-Meta Constraints.
\newblock Zenodo, 2026.
\newblock doi:10.5281/zenodo.18436828.

\bibitem{Takahashi2025FSNA}
K.~Takahashi.
\newblock Floor-Specified No-Meta Autonomy.
\newblock Zenodo, 2026.
\newblock doi:10.5281/zenodo.18258262.

\bibitem{Takahashi2025FMIC}
K.~Takahashi.
\newblock Friction-Minimal Intersubjective Contracts for No-Meta Autonomous Agents.
\newblock Zenodo, 2026.
\newblock doi:10.5281/zenodo.18308030.

\bibitem{Takahashi2025CLI}
K.~Takahashi.
\newblock Causal Loop Integrity for Observable-Only Backcasting: Strong No-Meta Supplement (Operational Core).
\newblock Zenodo, 2026.
\newblock doi:10.5281/zenodo.18335211.

\bibitem{RFC2119}
S.~Bradner.
\newblock Key words for use in RFCs to Indicate Requirement Levels.
\newblock RFC 2119, IETF, March 1997.
\newblock doi:10.17487/RFC2119.

\bibitem{RFC8174}
B.~Leiba.
\newblock Ambiguity of Uppercase vs Lowercase in RFC 2119 Key Words.
\newblock RFC 8174, IETF, May 2017.
\newblock doi:10.17487/RFC8174.

\bibitem{RFC4648}
S.~Josefsson.
\newblock The Base16, Base32, and Base64 Data Encodings.
\newblock RFC 4648, IETF, October 2006.
\newblock doi:10.17487/RFC4648.

\bibitem{RFC8785}
A.~Rundgren, B.~Jordan, and S.~Erdtman.
\newblock JSON Canonicalization Scheme (JCS).
\newblock RFC 8785, Independent Stream, June 2020.
\newblock doi:10.17487/RFC8785.

\bibitem{RFC9162}
B.~Laurie, E.~Messeri, and R.~Stradling.
\newblock Certificate Transparency Version 2.0.
\newblock RFC 9162, IETF, December 2021.
\newblock doi:10.17487/RFC9162.

\bibitem{FIPS1804}
Q.~H. Dang.
\newblock Secure Hash Standard (SHS).
\newblock FIPS PUB 180-4, National Institute of Standards and Technology (NIST), August 2015.
\newblock doi:10.6028/NIST.FIPS.180-4.

\bibitem{Merkle1988}
Ralph~C. Merkle.
\newblock A Digital Signature Based on a Conventional Encryption Function.
\newblock In {\em Advances in Cryptology -- CRYPTO '87}, Lecture Notes in Computer Science, vol. 293, pages 369--378. Springer, 1988.
\newblock doi:10.1007/3-540-48184-2\_32.

\bibitem{Bunz2020PCD}
Benedikt B{\"u}nz, Benjamin Fisch, Robert McNally, and Pratyush Mishra.
\newblock Proof-Carrying Data from Accumulation Schemes.
\newblock IACR Cryptology ePrint Archive, Report 2020/499, 2020.

\bibitem{Necula1997PCC}
George~C. Necula.
\newblock Proof-carrying code.
\newblock In {\em Proceedings of the 24th ACM SIGPLAN-SIGACT Symposium on
  Principles of Programming Languages (POPL)}, 1997.

\bibitem{Necula2000TV}
George~C. Necula.
\newblock Translation Validation for an Optimizing Compiler.
\newblock In {\em Proceedings of the ACM SIGPLAN Conference on Programming Language Design and Implementation (PLDI)}, 2000.
\newblock doi:10.1145/349299.349314.

\bibitem{Leroy2009CompCert}
Xavier Leroy.
\newblock Formal Verification of a Realistic Compiler.
\newblock {\em Communications of the ACM}, 52(7):107--115, 2009.
\newblock doi:10.1145/1538788.1538814.

\bibitem{BenSasson2014SNARK}
Eli Ben-Sasson, Alessandro Chiesa, Daniel Genkin, Eran Tromer, and Madars Virza.
\newblock {SNARKs for C}.
\newblock In {\em Advances in Cryptology -- CRYPTO}, 2014.

\bibitem{Groth2016}
Jens Groth.
\newblock On the Size of Pairing-based Non-interactive Arguments.
\newblock IACR Cryptology ePrint Archive, Report 2016/260, 2016.

\bibitem{BenSasson2018STARK}
Eli Ben-Sasson, Iddo Bentov, Yinon Horesh, and Michael Riabzev.
\newblock Scalable, Transparent, and Post-Quantum Secure Computational Integrity.
\newblock IACR Cryptology ePrint Archive, Report 2018/046, 2018.

\bibitem{Setty2021Nova}
Srinath Setty.
\newblock {Nova}: Recursive zero-knowledge arguments from folding schemes.
\newblock In {\em Advances in Cryptology -- CRYPTO}, 2021.

\end{thebibliography}

\end{document}
